\documentclass{article}
\usepackage[margin=1in]{geometry}

\setcounter{secnumdepth}{3}
% Useful packages
\usepackage{graphicx}
\usepackage{amsmath}
\usepackage{amssymb}
\usepackage{hyperref}
\usepackage{braket}
\usepackage{dsfont}
%\usepackage{showlabels}
\usepackage{amsthm}
\newtheorem{theorem}{Theorem}
\newtheorem{lemma}[theorem]{Lemma}
\newtheorem{corollary}[theorem]{Corollary}
\newtheorem{proposition}[theorem]{Proposition}
\theoremstyle{definition}
\newtheorem{definition}{Definition}[section]
\newtheorem{prop}{Conjecture}
\newtheorem{cor}{Corollary}
%-------------------packages-------------------------
\usepackage{enumerate}
\usepackage[utf8]{inputenc}
\usepackage{amsfonts}
\usepackage{amsbsy, mathtools}
% \usepackage{braket} % already loaded above
% \usepackage{graphicx} % already loaded above
\usepackage{xcolor}
%\usepackage{multirow}
% \usepackage{hyperref} % already loaded above
%\usepackage{caption}
%\usepackage{subcaption}
%\captionsetup{compatibility=false}
\usepackage{natbib}
%\usepackage{dcolumn}% Align table columns on decimal point
\usepackage{bm}% bold math
\usepackage{soul}
\usepackage{showlabels}

%Includes "References" in the table of contents
\usepackage[nottoc]{tocbibind}

\usepackage{bbm}
%\newtheorem{theorem}{Theorem}
%\newtheorem{lemma}[theorem]{Lemma}
%\newtheorem{corollary}[theorem]{Corollary}
%\theoremstyle{definition}
%\newtheorem{definition}{Definition}[section]
\usepackage{tikz}
\usepackage{quantikz}
\usepackage{textcomp}
\usepackage{caption}
\usepackage{microtype}
\newcommand{\Tr}{\text{tr}}
\newcommand{\mA}{\mathcal{A}}
\newcommand{\mB}{\mathcal{B}}
\newcommand{\nn}{\nonumber}
\newcommand{\mL}{\mathcal{L}}
\newcommand{\lb}{\left(}
\newcommand{\rb}{\right)}
\newcommand{\mU}{\mathcal{U}}
\newcommand{\mH}{\mathcal{H}}
\usepackage[T1]{fontenc}
\usepackage{lmodern}
\usepackage{subcaption}
\newcommand{\tOO}{\tilde{\mathcal{O}}}
\newcommand{\mO}{\mathcal{O}}
\newcommand{\mK}{\mathcal{K} }
\newcommand{\ep}{\epsilon}
\newcommand{\mbR}{\mathbb{R}}
\newcommand{\mM}{\mathcal{M}}
\newcommand{\mT}{\mathcal{T}}
\newcommand{\NL}[1]{\textcolor{blue}{#1}}
\newcommand{\p}{\partial}
\newcommand{\tr}{\text{tr}}
\newcommand{\tT}{\tilde{T}}

\usepackage{graphicx} % Required for inserting images

\title{Information loss and ensemble-averaging theories}
\author{Nima Lashkari}
\date{January 2026}

\begin{document}
    \maketitle

    \section{Results}



    We find that in time-independent ensemble averaged quantum theories with a large
    $N$ limit, the $k$-twirl map becomes
    \begin{eqnarray}
        \mathcal{E}_g\lb U^{\otimes k}\otimes (U^*)^{\otimes k} \rb=e^{-t^2\mathcal{G}_{eff}^{(k)}}
    \end{eqnarray}
    in the case of the infinite temperature state. At finite temperature, the
    result will be \NL{finish}

    In theories with Brownian random couplings, since the Hamiltonian is time-dependent,
    there is no canonical notion of thermal state. Instead, we will focus on the
    infinite temperature state. In this case, we can perform the average exactly
    for all $N$, and the result is
    \begin{eqnarray}
        \mathcal{E}_g\lb U^{\otimes k}\otimes (U^*)^{\otimes k} \rb=e^{-t\mathcal{G}_{eff}^{(k)}}
    \end{eqnarray}
    with the same generator $\mathcal{G}_{eff}^{(k)}$.

    We analyze the spectral properties of $\mathcal{G}_{eff}^{(k)}$ and find that
    \NL{finish}

    We comment on the $q$-deformed limit $N\to \infty$ with $q^{2}/N$ kept fixed.

    \section{To do list:}
    \begin{enumerate}
        \item Finish the generalize the derviation of the $k$-twirl map in regular
            SYK to finite temperature at large $N$.

        \item Understand the physical implications of the $t^{2}$ in the exponent
            of the $k$-twirl map in the large $N$ limit of regular SYK. At late
            time, the decay seems too fast.

        \item Establish the spectral gap in the spectrum of the generator $\mathcal{G}
            _{eff}^{(k)}$.

        \item Think about the gluing assumption 1.4 in Nick's notes. Does it
            hold in SYK? Related issue is this: if we couple two Brownian SYK
            models as Swingle does in his 2021 paper, does it look like the Alice-Bob
            protocol for achieving a $k$-twirl map on a combined system $AB$ by applying
            the $k$-twirl maps on $A$ and $B$, separately, swapping $r$ qubits,
            and applying the local $k$-twirls again?
    \end{enumerate}

    \section{Correlation functions and $k$-twirl maps}

    Let $\mH$ be the microscopic Hilbert space (dimension $d$) and $\mA=\mathrm{End}
    (\mH)$ the operator algebra. We use the doubled (Choi/GNS) Hilbert space
    \begin{equation}
        \mH_{\GNS}:=\mH_{L}\otimes \mH_{R}, \qquad \ket{1}:=\frac{1}{\sqrt{d}}\sum
        _{i=1}^{d}\ket{i}_{L}\ket{i}_{R}.
    \end{equation}
    We identify an operator $A\in\mA$ with the vector
    $\ket{A}:=(A_{L}\otimes 1_{R})\ket{1}$. With this choice, for all
    $A,B\in\mA$,
    \begin{equation}
        \braket{1|A_L\otimes B_R|1}=\frac{1}{d}\tr\big(A\,B^{\mathsf{T}}\big),
    \end{equation}
    where $\mathsf{T}$ denotes transpose in the basis used to define $\ket{1}$.

    \subsection*{Permutation action and cycle states}

    Consider $k$ copies of the GNS Hilbert space, $\mH_{\GNS}^{\otimes k}$, with
    canonical vector $\ket{1^{\otimes k}}$. We choose the representation of the permutation
    group $S_{k}$ which acts on the left factors:
    \begin{equation}
        \mathbb{S}_{\pi}\Big(\bigotimes_{a=1}^{k}\ket{i_a}_{L}\ket{j_a}_{R}\Big)
        \,=\,\bigotimes_{a=1}^{k}\ket{i_{\pi(a)}}_{L}\ket{j_a}_{R},\qquad \pi\in
        S_{k}.
    \end{equation}
    For a $k$-cycle $\pi\in C_{k}\subset S_{k}$ and any operators
    $\mO_{1},\dots, \mO_{k}\in\mA$ one has the standard identity
    \begin{equation}
        \bra{\mathbb{S}_\pi 1^{\otimes k}}\Big(\prod_{a=1}^{k}(\mO_{a})_{L}^{(a)}
        \Big)\ket{1^{\otimes k}}\,=\,\tr\big(\mO_{\pi(k)}\cdots \mO_{\pi(1)}\big)
        .
    \end{equation}
    In particular, for the cyclic shift $\pi_{0}=(1\,2\,\dots\,k)$ this gives a cyclic
    trace.

    \subsection*{Complexified evolution as a superoperator}

    We work with dimensionless variables (set $\beta=1$), so that the thermal state
    is $\rho\equiv\rho_{\beta=1}=e^{-H}/Z$ with $Z=\mathrm{Tr}(e^{-H})$.
    Introduce a complexified time
    \begin{equation}
        z=\theta+i t,\qquad \theta\in[0,\tfrac12],\quad t\in\mathbb{R},
    \end{equation}
    and define
    \begin{equation}
        U_{z}:=e^{-zH}.
    \end{equation}
    The corresponding completely-positive ``thermalized'' adjoint action is
    \begin{equation}
        \boxed{\;\mathcal{U}_{z}(\mO):=U_{z}\,\mO\,U_{z}^{\dagger} =e^{-(\theta+i t)H}\,\mO\,e^{-(\theta-i t)H}.\;}
    \end{equation}
    (Note that $\theta\ge 0$ produces thermal damping rather than growth.)

    On the doubled space, left- and right-multiplication are represented by
    \begin{equation}
        \mathbb{U}_{z}:=U_{z,L}\otimes \big(U_{z}^{\dagger}\big)_{R}^{\mathsf{T}}
        =e^{-zH_L}\otimes e^{-z^{\ast} H_R^{\mathsf{T}}}.
    \end{equation}
    For $\vec z=(z_{1},\dots,z_{k})$ define
    \begin{equation}
        \mathbb{U}_{\vec z}:=\bigotimes_{a=1}^{k}\mathbb{U}_{z_a}^{(a)}.
    \end{equation}
    Then the trace of a product of superoperators has the GNS/cycle
    representation
    \begin{equation}
        \boxed{\; \tr\big(\mathcal{U}_{z_1}(\mO_1)\cdots \mathcal{U}_{z_k}(\mO_k)\big) \,=\,\bra{\mathbb{S}_{\pi_0}1^{\otimes k}}\,\mathbb{U}_{\vec z}\,\bigotimes_{a=1}^{k}(\mO_a)_{L}\ket{1^{\otimes k}}.\;}
    \end{equation}
    Thermal correlators arise by imposing the ``thermal strip'' constraint
    \begin{equation}
        \boxed{\;\sum_{a=1}^k \Re z_a=\tfrac12\;}
    \end{equation}
    (which guarantees $\tr(\mathcal{U}_{z_1}(1)\cdots \mathcal{U}_{z_k}(1))=\tr (
    e^{-H})=Z/d$).

    \subsection*{Thermal one- and two-point functions}

    \paragraph{One-point.}
    For any operator $\mO$ and any $t\in\mathbb{R}$,
    \begin{equation}
        \label{1ptthermal}\boxed{\; \tr(\rho\,\mO) =\frac{\bra{1}\,\mathbb{U}_{\frac{1}{2}+i t}\,(\mO)_{L}\ket{1}}{\bra{1}\,\mathbb{U}_{\frac{1}{2}+i t}\ket{1}} =\frac{\tr\big(e^{-\frac{1}{2}H}e^{-i tH}\mO\,e^{+i tH}e^{-\frac{1}{2}H}\big)}{\tr(e^{-H})}=\frac{\tr\big(e^{-H}\mO\big)}{\tr(e^{-H})}\;}
    \end{equation}
    The ratio is $t$-independent by cyclicity of the trace.

    \paragraph{Two-point (Wightman/KMS form).}
    Fix $\tau\in[0,\tfrac12]$ and define the thermal Wightman function
    \begin{equation}
        \label{2ptthermal}G_{12}(t,\tau):=\frac{\tr\big(e^{-(1-\tau)H}\,\mO_{1}(t)\,e^{-\tau
        H}\,\mO_{2}\big)}{\tr(e^{-H})}, \qquad \mO_{1}(t):=e^{+iHt}\mO_{1}e^{-iHt}
        .
    \end{equation}
    It can be written using $k=3$ copies (with the third insertion taken to be the
    identity) as
    \begin{equation}
        \boxed{\; G_{12}(t,\tau) =\frac{\bra{\mathbb{S}_{\pi_0}1^{\otimes 3}}\,\mathbb{U}_{(z_1,z_2,z_3)}\,\big((\mO_{1})_{L}\otimes (\mO_{2})_{L}\otimes 1_{L}\big)\ket{1^{\otimes 3}}}{\bra{\mathbb{S}_{\pi_0}1^{\otimes 3}}\,\mathbb{U}_{(z_1,z_2,z_3)}\ket{1^{\otimes 3}}} ,}
    \end{equation}
    with the convenient choice
    \begin{equation}
        (z_{1},z_{2},z_{3})=(0,\,\tau,\,\tfrac12-\tau).
    \end{equation}
    (Real-time dependence may equivalently be kept in $\mO_{1}(t)$ as above, or
    absorbed into $z_{2}\to \tau+i t$.)

    \subsection*{General $(k\! -\! 1)$-point functions}

    For $k\ge 2$ we define a thermal $(k-1)$-point function by inserting the
    ideneity operator $1$ as the $k$-th operator:
    \begin{equation}
        \boxed{\; G_{1\cdots (k-1)}(\vec z) :=\frac{\bra{\mathbb{S}_{\pi_0}1^{\otimes k}}\,\mathbb{U}_{\vec z}\,\big((\mO_{1})_{L}\otimes\cdots\otimes (\mO_{k-1})_{L}\otimes 1_{L}\big)\ket{1^{\otimes k}}}{\bra{\mathbb{S}_{\pi_0}1^{\otimes k}}\,\mathbb{U}_{\vec z}\ket{1^{\otimes k}}},\qquad \sum_{a=1}^k\Re z_a=\tfrac12.\;}
    \end{equation}
    A concrete parametrization adapted to Euclidean insertion times $\tau_{1},\cdots
    \tau_{k-1}$ with $0=\Re\tau_{1}\le \Re\tau_{2}\le\cdots\le \Re\tau_{k-1}\le \tfrac
    12$ (which fixes the overall time-translation along the thermal circle) is
    \begin{align}
         & G_{1\cdots (k-1)}(\tau_{1},\cdots, \tau_{k-1})\equiv\braket{\mO_1(\tau_1)\cdots \mO_{k-1}(\tau_{k-1})}=G_{1\cdots(k-1)}(\vec{z})\nn \\
         & \Re\tau_{1}=\Re z_{1},\qquad \Re\tau_{k}=1/2-\Re\tau_{k-1}\nn                                                                       \\
         & \forall 1<j<k:\qquad \Re\tau_{j}=\sum_{a=1}^{j}(\Re z_{a}+\Re z_{a+1}),\qquad \tau_{j}=\Re \tau_{j}+i\Im z_{j}
    \end{align}

    %\end{document}

    \section{Gaussian regulator from a time-independent ensemble}

    Consider a quantum system with Hamiltonian
    \begin{eqnarray}
        H(g)=\sum_I g_I V_I
    \end{eqnarray}
    where $V_{I}$ are dimensionless self-adjoint operators and $g_{I}$ are real
    coupling parameters with units of energy. We choose $g_{I}$ at random with a
    probability measure that, for simplicity, we take to be Gaussian:
    \begin{eqnarray}
        d\mu(g)= \prod_I\lb dg_I
        \frac{e^{-\frac{1}{2\sigma}\sum_I g_I^2}}{\sqrt{2\pi \sigma}} \rb
    \end{eqnarray}
    where $\sigma$ has units of energy squared and sets the width of the
    distribution, i.e., $[\sigma]=2$ and $[g_{I}]=1$. The first moments are
    \begin{eqnarray}
        \mathcal{E}_g(g_I)=0, \qquad \mathcal{E}_g(g_I g_J)=\delta_{IJ}\sigma\ .
    \end{eqnarray}

    We represent the thermal state $\rho_{\beta}\propto e^{-\beta H}$ by its canonical
    purification (thermofield double) $\ket{\rho_\beta^{1/2}}$ and the corresponding
    GNS Hilbert space $\mathcal{H}_{\beta}$, generated by the action of
    observables $\mO$ on this vector:
    \begin{eqnarray}
        \mO\to \mO\ket{\rho_\beta^{1/2}}\ .
    \end{eqnarray}
    It is convenient to choose $\beta$ to set the unit of time and define dimensionless
    parameters and Hamiltonian
    \begin{eqnarray}
        &&\tilde{g}_I=\beta g_I,\qquad \tilde{\sigma}=\beta^2 \sigma, \qquad
        \tilde{t}=t/\beta\nn\\
        &&\tilde{H}=\sum_I \tilde{g}_I V_I
    \end{eqnarray}
    The thermal state is a function of $\tilde{g}_{I}$; its canonical
    purification is
    \begin{eqnarray}
        \ket{\rho_\beta^{1/2}}=\ket{\rho_{\tilde{g}}^{1/2}}=(e^{-\tilde{H}/2}\otimes
        1)\ket{1}\ .
    \end{eqnarray}
    where
    \begin{eqnarray}
        \ket{1}\equiv \frac{1}{\sqrt{d}}\sum_E \ket{E E}
    \end{eqnarray}
    and $\{\ket{E}\}$ is the eigenbasis of the Hamiltonian. Since, in general, members
    of our ensemble of Hamiltonians do not commute, we choose some basis for
    $\ket{1}$, and then use the mirror equation
    \begin{eqnarray}
        (\mathcal{O}\otimes 1)\ket{1}=(1\otimes \mathcal{O}^T)\ket{1}
    \end{eqnarray}
    where the transpose is defined in the basis of $\ket{1}$. For self-adjoint
    operators like terms in the Hamiltonians, we have
    \begin{eqnarray}
        (V_I\otimes 1)\ket{1}=(1\otimes V_I^*)\ket{1}\ .
    \end{eqnarray}

    The measure for our ensemble average is
    \begin{eqnarray}
        d\mu(\tilde{g})=\prod_I\lb
        \frac{e^{-\tilde{g}_I^2/(2\tilde{\sigma})}}{\sqrt{2\pi\tilde{\sigma}}} d\tilde{g}_I\rb
    \end{eqnarray}
    For fixed $\tilde{g}$, the time evolution is a unitary super-operator
    \begin{eqnarray}
        \mO(t)=e^{-i H t}\mO e^{i H t}=e^{-i \tilde{H}\tilde{t}} \mO e^{i\tilde{H}\tilde{t}}\ .
    \end{eqnarray}
    It is represented on the Hilbert space as the GNS action of the modular Hamiltonian
    $\tilde{H}$:
    \begin{eqnarray}
        &&\mO(t)\ket{1}=e^{-i \tilde{H} \tilde{t}}\mO\ket{1}\nn\\
        &&\tilde{H}=\sum_I \tilde{g}_I \hat{V}_{I},\qquad \hat{V}_I=V_{I,L}-V^T_{I,R},\qquad
        \hat{V}_I\ket{1}=0.
    \end{eqnarray}
    Here, it is important that $[e^{i \tilde{H} \tilde{t}},e^{-\tilde{H}/2}]=0$.

    The $g$-averaged time evolution is a unital linear CP map
    \begin{eqnarray}
        \mathcal{E}_g(\mO(t))=\int d\mu(g) e^{-i H(g) t}\mO e^{i H(g) t}\ .
    \end{eqnarray}
    %which as we will show below does not preserve the thermal state $\rho_\beta$!
    In the GNS Hilbert space, the unital CP map $\mathcal{E}_{g}(\cdot)$
    corresponds to the action of the operator
    \begin{eqnarray}
        \mathcal{E}_g(\mO(t))\ket{1}=\int d\mu(g) e^{-i \tilde{t}\sum_I \tilde{g}_I\hat{V}_I }\mO\ket{1}\ .
    \end{eqnarray}

    \subsection{Commuting case}

    When there is only a single coupling $g$ (i.e. one index $I$), the Gaussian
    integral can be performed explicitly,
    \begin{eqnarray}
        \int \frac{dg}{\sqrt{2\pi \tilde{\sigma}}}\:e^{-\frac{g^{2}}{2\tilde{\sigma}}}\:e^{-i\tilde{t}g\hat{V}}
        =e^{-\frac{\tilde{\sigma}}{2}\tilde{t}^2\hat{V}^2}\ .
    \end{eqnarray}
    In fact, this integral can be extended to the complex $\tild{t}$ domain. Therefore,
    for the $k$-twirl map we have
    \begin{align}
        \int d\mu(g) \mathbb{U}_{\vec{z}}=e^{-\frac{\tilde{\sigma}}{2}\lb\sum_{a=1}^k (iz_a)\hat{V}_{(a)}\rb^2}
    \end{align}
    and
    \begin{align}
        G_{1\cdots (k-1)k}(\tau_{1},\cdots,\tau_{k}) & =\braket{\mathbb{S}_{\pi}1^{\otimes k}|e^{-\frac{\tilde{\sigma}}{2}\lb\sum_{a=1}^k (iz_a)\hat{V}_a\rb^2} \otimes_{a=1}^k (\mO_{a})_L|1^{\otimes k}}\nn                   \\
                                                     & = \braket{1^{\otimes k}|e^{\frac{\tilde{\sigma}}{2}\lb\sum_{a=1}^k (iz_a)(V_L^{\pi(a)}-(V_R^{(a)})^T)\rb^2} \mathbb{S}_\pi \otimes_{a=1}^k (\mO_{a})_L|1^{\otimes k}}\nn \\
                                                     & = \braket{1^{\otimes k}|e^{\frac{\tilde{\sigma}}{2}\lb\sum_{a=1}^k (iz_a)(V_L^{\pi(a)}-(V_R^{(a)}^T))\rb^2} \otimes_{a=1}^k (\mO_{\pi(a)})_L|1^{\otimes k}}
    \end{align}
    where we have used
    \begin{align}
        \mathbb{S}_{\pi}(\otimes_{a=1}^{k}(\mO_{a})_{L}) \mathbb{S}_{\pi}^{\dagger}= \otimes_{a=1}^{k}(\mO_{\pi(a)})_{L}.
    \end{align}
    We also have
    \begin{align}
        (V_{L}^{(\pi(a))}-(V_{R}^{(a)})^{T})\ket{1^{\otimes k}}= (V_{L}^{(\pi(a))}-V_{L}^{(a)})\ket{1^{\otimes k}}
    \end{align}
    Therefore, we can write
    \begin{align}
        G_{1\cdots (k-1)k}(\tau_{1},\cdots,\tau_{k}) & =\braket{1^{\otimes k}|e^{\frac{\tilde{\sigma}}{2}\lb\sum_{a=1}^k (iz_a)(V_L^{\pi(a)}-V_L^{(a)})\rb^2} \otimes_{a=1}^k (\mO_{\pi(a)})_L|1^{\otimes k}}
    \end{align}

    It is clear from (\ref{1ptthermal})that for the one-point function individual
    elements of the ensemble is time-independent. For the $g$-averaged thermal
    two-point function with insertions at Euclidean times $\tau_{1}$ and
    $\tau_{2}$ we have
    \begin{align}
        G_{12}(\tau_{1},\tau_{2}) & =\braket{\mathbb{S}_{\pi_0}1^{\otimes 3}| e^{-\frac{\tilde{\sigma}}{2}\lb\tau\hat{V}_2+(1/2-\tau)\hat{V}_3\rb^2} ((\mO_1)_L\otimes (\mO_2)_L\otimes 1_L)|1^{\otimes 3}} %%
    \end{align}

    % therefore, we have
    % \begin{align}
    %     \braket{\mathbb{S}_{\pi}1^{\otimes 3}|e^{-\frac{\tilde{\sigma}}{2}\lb\tau\hat{V}_2+(1/2-\tau)\hat{V}_3\rb^2} ((\mO_1)_L\otimes (\mO_2)_L\otimes 1_L)|1^{\otimes 3}} & = \braket{1^{\otimes 3}|e^{\frac{\tilde{\sigma}}{2}\lb\tau\hat{V}_2+(1/2-\tau)\hat{V}_3\rb^2} \mathb{S}_\pi ((\mO_1)_L\otimes (\mO_2)_L\otimes 1_L)|1^{\otimes 3}}\nn \\
    %                                                                                                                                                                         & = \braket{1^{\otimes 3}|e^{\frac{\tilde{\sigma}}{2}\lb\tau\hat{V}_2+(1/2-\tau)\hat{V}_3\rb^2} ((\mO_1)_L\otimes (\mO_2)_L\otimes 1_L)|1^{\otimes 3}}
    % \end{align}

    % The $g$-averaged thermal one-point function is
    % \begin{eqnarray}
    %     \braket{\rho_\beta^{1/2}|\mathcal{E}_g\mO(t)\rho_\beta^{1/2}}&=\braket{e^{-\frac{\tilde{\sigma}}{2}\tilde{t}^2\hat{V}^2}\rho_\beta^{1/2}|\mO\rho_\beta^{1/2}}\nn\\
    %     &=
    % \end{eqnarray}
    % which is clearly time-dependent. Hence, $\mathcal{E}_{g}$ does not preserve
    % $\rho_{\beta}$. In particular, in the limit $\tilde{t}\to \infty$, the map projects
    % onto the kernel of $\hat{V}^{2}$. The $g$-averaged two-point function is
    % \begin{eqnarray}
    %     &&\mathcal{E}_g(\braket{\mO(t_1)\rho_\beta^{1/2}|\mO(t_2)\rho_\beta^{1/2}})=\braket{e^{-\frac{\tilde{\sigma}}{2}\hat{V}^2(t_1-t_2)^2}\mO\rho_\beta^{1/2}|\mO\rho_\beta^{1/2}}
    % \end{eqnarray}
    % which is translation-invariant.
    % \NL{When there is only one $g$, for the eigenvectors of $\tilde{V}$ of frequency $\omega$:
    %   \begin{eqnarray}
    %     &&e^{-i g V t}(\mO_\omega)e^{i g V t}=e^{-i\omega g t}\mO_\omega\nn\\
    %     &&\mathcal{E}_g\tr(\rho_\beta \mO_\omega(t))=e^{-\frac{\sigma}{2}(\omega t)^2}\tr(\rho_\beta \mO_\omega)\ .
    %   \end{eqnarray}
    % However, note that these are not modular eigenmodes, because the modular Hamiltonian is $\tilde{H}=g \tilde{V}$ and not just $\tilde{V}$.}

    % \NL{For a single $g$, note that if instead of $g$ we average over modular time with measure
    %   \begin{eqnarray}
    %     d\nu(t)=\sqrt{\frac{\sigma}{2\pi}}e^{-\frac{\sigma t^2}{2}}dt
    %   \end{eqnarray}
    %   we find
    %   \begin{eqnarray}
    %     \int d\nu(t) e^{-i g \tilde{V}t}\mO\ket{\rho^{1/2}}=e^{-\frac{g^2}{2\sigma}\tilde{V}^2}\mO\ket{\rho_\beta^{1/2}}\ .
    %   \end{eqnarray}
    % Does this imply a relation between $g$-average and time average? Should we have started with $H=V$ from the start?}

    More generally, if all the operators $\hat{V}_{I}$ mutually commute, the
    same calculation can be applied in a common eigenbasis and yields
    \begin{eqnarray}
        \int d\mu(g)\:e^{-i\tilde{t}\sum_I g_I\hat{V}_I}=e^{-\frac{\tilde{\sigma}}{2}\tilde{t}^2\sum_I \hat{V}_I^2}\ .
    \end{eqnarray}
    Physically, this acts as a Gaussian regulator. This also extends to the case
    where

    For noncommuting $\hat{V}_{I}$, the Gaussian integral above does not
    collapse to a single exponential in $\sum_{I}\hat{V}_{I}^{2}$. An exact expression
    can be written as an ordered Wick expansion. Expanding
    $e^{-i\tilde{t}\sum_I g_I\hat{V}_I}$ in a power series and using Wick's
    theorem gives
    \begin{eqnarray}
        &&\int d\mu(g)\:e^{-i\tilde{t}\sum_I g_I\hat{V}_I}\nn\\
        &&\qquad=\sum_{n=0}^\infty \frac{(-1)^{n}}{(2n)!}\tilde{t}^{2n}\sum_{I_1,\ldots,I_{2n}}\mathcal{E}_g\lb
        g_{I_1}\cdots g_{I_{2n}}\rb\;\hat{V}_{I_1}\cdots\hat{V}_{I_{2n}}\nn\\
        &&\qquad=\sum_{n=0}^\infty \frac{(-1)^{n}}{2^{n}n!}\tilde{\sigma}^n\tilde{t}^{2n}\sum_{\text{pairings }P}\sum_{I_1,\ldots,I_n}\hat{V}_{I_{\pi_P(1)}}\hat{V}_{I_{\pi_P(2)}}\cdots\hat{V}_{I_{\pi_P(2n)}}\ ,
    \end{eqnarray}
    where $P$ runs over pairings of $\{1,\ldots,2n\}$ and $\pi_{P}$ is the map
    that repeats each label $I_{r}$ in the two paired slots. The operator ordering
    is inherited from the product inside $\hat{V}_{I_1}\cdots\hat{V}_{I_{2n}}$.
    In particular,
    \begin{eqnarray}
        \int d\mu(g)\:e^{-i\tilde{t}\sum_I g_I\hat{V}_I} &&=1-\frac{\tilde{\sigma}}{2}\tilde{t}^2\sum_I
        \hat{V}_I^2\nn\\
        &&\qquad+\frac{\tilde{\sigma}^{2}}{8}\tilde{t}^4\sum_{I,J}\lb \hat{V}_I^2\hat{V}_J^2+\hat{V}_I\hat{V}_J\hat{V}_I\hat{V}_J+\hat{V}_I\hat{V}_J^2\hat{V}_I\rb+O(\tilde{t}^6)\ ,
    \end{eqnarray}
    Eq.~(17) can also be written as an exponential, but in general it is the exponential
    of an infinite series of nested commutators (a Magnus type expansion).
    Concretely, one may write
    \begin{eqnarray}
        \int d\mu(g)\:e^{-i\tilde{t}\sum_I g_I\hat{V}_I}=\exp\big(\Omega(\tilde{t})\big),\qquad
        \Omega(\tilde{t})=\sum_{n\geq 1}\Omega_n(\tilde{t})\ .
    \end{eqnarray}
    where each $\Omega_{n}$ is built from $n$-fold commutators of the $\hat{V}_{I}$
    and starts at order $\tilde{t}^{2n}$ (odd orders vanish by Gaussianity). In
    general, the $\hat{V}_{I}$ do not commute, $\Omega(\tilde{t})$ does not
    truncate, and we can write the full integral as
    \begin{eqnarray}
        \int d\mu(g)\:e^{-i\tilde{t}\sum_I g_I\hat{V}_I}=\lb \lb e^{\frac{\tilde{\sigma}^{2}}{2}\sum_I \frac{\p^{2}}{\p g^{2}_{I}}}\rb
        e^{-i \tilde{t}\sum_I g_I \hat{V}_I}\rb_{g_I=0}\ .
    \end{eqnarray}

    \paragraph{The $k$-twirl channel and its moment operator.}
    For a fixed realization $J$, define $U_{J}(t)=e^{-itH(J)}$ and the $k$-fold
    adjoint action
    \begin{equation}
        \mathcal{U}^{(k)}_{J,t}(Y)\;:=\;U_{J}(t)^{\otimes k}\,Y\,U_{J}(t)^{\dagger\otimes
        k}, \qquad Y\in \mathrm{End}(\mathcal{H}^{\otimes k}).
    \end{equation}
    The disorder-averaged $k$-twirl map is
    \begin{equation}
        \mathcal{E}^{(k)}_{t}(Y)\;:=\;\mathbb{E}_{J}\!\left[\mathcal{U}^{(k)}_{J,t}
        (Y)\right].
    \end{equation}

    Introduce the vectorization/Choi-GNS map $Y\mapsto |Y\rangle\rangle$ on a doubled
    space,
    \begin{equation}
        |Y\rangle\rangle := (Y\otimes \mathbf{1})\,|1\rangle, \qquad |1\rangle :=
        \frac{1}{\sqrt{d}}\sum_{E}|E\rangle\otimes |E\rangle,
    \end{equation}
    so that left- and right-actions are represented by
    \begin{equation}
        (O_{L}- O_{R}^{T})\,|Y\rangle\rangle = |[O,Y]\rangle\rangle.
    \end{equation}
    Then the averaged moment operator (Choi operator) is
    \begin{equation}
        M_{t}^{(k)}\;:=\;\mathbb{E}_{J}\!\left[U_{J}(t)^{\otimes k}\otimes \big(U
        _{J}(t)^{*}\big)^{\otimes k}\right], \qquad |\mathcal{E}_{t}^{(k)}(Y)\rangle
        \rangle = M_{t}^{(k)}\,|Y\rangle\rangle.
    \end{equation}

    %------------------------------------------------------------
    \paragraph{Hatted operators and Gaussian integral.}
    Define, for each replica $a=1,\dots,k$,
    \begin{equation}
        \widehat H_{J}^{(a)}:= H_{J,L}^{(a)}- \big(H_{J,R}^{(a)}\big)^{T}, \qquad
        \widehat H_{J}^{(k)}:= \sum_{a=1}^{k}\widehat H_{J}^{(a)}.
    \end{equation}
    Since $H(J)=\sum_{I}J_{I}V_{I}$, we can write
    \begin{equation}
        \widehat H_{J}^{(k)}\;=\; \sum_{I}J_{I}\,\widehat V_{I}^{(k)}, \qquad \widehat
        V_{I}^{(k)}:= \sum_{a=1}^{k}\widehat V_{I}^{(a)}, \qquad \widehat V_{I}^{(a)}
        := V_{I,L}^{(a)}- \big(V_{I,R}^{(a)}\big)^{T}.
    \end{equation}
    Therefore
    \begin{equation}
        M_{t}^{(k)}\;=\; \mathbb{E}_{J}\!\left[e^{-it\,\widehat H_J^{(k)}}\right]
        \;=\;\int d\mu(J)\;\exp\!\left(-it\sum_{I}J_{I}\,\widehat V_{I}^{(k)}\right
        ), \label{eq:Mt_as_gaussian_integral}
    \end{equation}
    where $d\mu(J)$ is the product Gaussian measure.

    More generally, we can consider
    \begin{eqnarray}
        M(\tau_1,\cdots,\tau_k)&&:=\mathbb{E_g}\!\left[\otimes_{a=1}^k \lb e^{-\tau_a H_L^{(a)})}\otimes
        \big(e^{-\tau_a^* H_R^{(a)}}\big)\rb^{\otimes k}\right],\nn\\
        &&=\mathbb{E}_g\left[ e^{-i \sum_I g_I \mathbb{F}_I}\right]\nn\\
        \mathbb{F}_I&&=\sum_{a=1}^k(\tau_a V^{(a)}_{I,L}-\tau_a^* V^{(a)}_{I,R})\nn\\
        &&\tau_a=\beta_a/4+it_a\ .
    \end{eqnarray}
    Note that as a superoperator the operator above for $k=1$ acts as
    \begin{eqnarray}
        &&\mathcal{F}(Y)=e^{-(\beta/4+it)}Y e^{-(\beta/4-it)}\nn\\
        &&(\mathcal{F}(Y_L)\otimes 1_R)\ket{e}=(U_t Y_L U_t^\dagger\otimes 1_R)\ket{\rho_\beta^{1/2}}\ .
    \end{eqnarray}

    %------------------------------------------------------------
    \paragraph{Ordered Wick expansion and leading Gaussian regulator.}
    Expanding the exponential and averaging with Wick's theorem gives the exact ordered
    pairing series
    \begin{equation}
        M_{t}^{(k)}= \sum_{n=0}^{\infty}\frac{(-1)^{n}}{2^{n}n!}\,(\sigma t^{2})^{n}
        \sum_{\text{pairings }P}\; \sum_{I_1,\dots,I_n}\widehat V_{I_{\pi_P(1)}}^{(k)}
        \widehat V_{I_{\pi_P(2)}}^{(k)}\cdots \widehat V_{I_{\pi_P(2n)}}^{(k)}, \label{eq:ordered_wick_quenched}
    \end{equation}
    so in particular
    \begin{equation}
        M_{t}^{(k)}= \mathbf{1}-\frac{\sigma t^{2}}{2}\sum_{I}\big(\widehat V_{I}
        ^{(k)}\big)^{2}+\mathcal{O}(t^{4}). \label{eq:Mt_smallt}
    \end{equation}

    Equivalently, one may write a Magnus/cumulant form
    \begin{equation}
        M_{t}^{(k)}=\exp\!\big(\Omega(t)\big), \qquad \Omega(t)=\Omega_{1}(t)+\Omega
        _{2}(t)+\cdots,
    \end{equation}
    where $\Omega_{1}(t)=-(\sigma t^{2}/2)\sum_{I}(\widehat V_{I}^{(k)})^{2}$,
    and $\Omega_{n\ge 2}(t)$ are built from nested commutators and begin at
    order $t^{2n}$.

    \subsection{Gaussian regulator from large $N$ limit}

    Now, consider the system to be comprised of $N$ degrees of freedom, $i=1,\cdots
    , N$, and $V_{I}$ is an interaction term that couples different sites. For instance,
    if the sites sit on a lattice, it is natural to consider nearest neighbor two-body
    interactions $V_{I}$ with $I=\{i,j\}$ that are adjacent on the lattice. More
    generally, we can consider the interaction to be non-local and $q$-body as in
    the case of hypergraphs such that the label $I$ is a set of $q$ distinct sites
    $I=\{i_{1},i_{2},\cdots i_{q}\}$. Assuming that the order of labels does not
    matter have the interaction
    \begin{eqnarray}
        &&H=q!\sum_{i_1<\cdots <i_q}g_{i_1\cdots i_q} V_{i_1\cdots i_q}\nn\\
        &&V_{i_1\cdots i_q}=\mathcal{O}_{i_1}\cdots \mathcal{O}_{i_q}\ .
    \end{eqnarray}
    If the operators $[\mathcal{O}_{i},\mathcal{O}_{j}]\propto \delta_{ij}$, then
    $[V_{I},V_{J}]=0$ for all $I,J$, and we obtain
    \begin{eqnarray}
        \int d\mu(g)e^{i t \sum_I g_I \hat{V}_I}=e^{-\sigma t^2\sum_I \hat{V}_I^2/2}\ .
    \end{eqnarray}
    If there are Majorana fermions sitting at each site, i.e.
    $\{\psi_{i},\psi_{j}\}=2\delta_{ij}$, there are certain non-vanishing $[V_{I}
    ,V_{J}]$. If the two sets of indices $I$ and $J$ have no overlaps, i.e.
    $I\cap J=\emptyset$, we have $[V_{I},V_{J}]=0$. We will assume that $q$ is
    even, then if the number of overlaps between the two sets $m=|I\cap J|$ is even,
    the potential terms commute, and if it is odd, they anticommute:
    \begin{eqnarray}
        [V_I,V_J]=(1-(-1)^{|I\cap J|})V_IV_J\ .
    \end{eqnarray}
    The case with a Majorana fermion at each siste is the SYK model that we will
    discuss in the next section. Here, we simply make the following observation.
    Choosing random sets $I,J$ the probability of having at least one index in
    common is
    \begin{eqnarray}
        P(\text{overlap})=1-\frac{{N-q\choose q}}{{N \choose q}}
    \end{eqnarray}
    which in the limit of fixed $q$ and large $N$ scales as $q^{2}/N\to 0$. Therefore,
    we can ignore the commutators, and obtain the Gaussian regulator map
    \begin{eqnarray}
        \int d\mu(g) e^{-t\sum_I g_I \hat{V}_I}=e^{-\sigma t^2\sum_I \hat{V}_I^2/2}\ .
    \end{eqnarray}
    Below, we will see that this large $N$ behavior generalizes to the $k$-twirl
    map, as well.

    \subsection{Gaussian regulator for the $k$-twirl at large $N$}
    \label{subsec:quenched_syk_gaussian_regulator}
    %------------------------------------------------------------
    \NL{Decide whether to use $g_{I}$ or $J_{I}$ as coupling}

    Let $H(J)$ be a (quenched) SYK Hamiltonian written as a Gaussian linear combination
    \begin{equation}
        H(J)\;=\; i^{q/2}\sum_{I}J_{I}\,V_{I},
        %  \qquad
        % V_{i_1<\cdots<i_q} \;=\; \chi_{i_1}\chi_{i_2}\cdots \chi_{i_q},
    \end{equation}
    with independent centered Gaussian couplings
    \begin{equation}
        \mathbb{E}[J_{I}]=0, \qquad \mathbb{E}[J_{I}J_{J}]=\sigma\,\delta_{IJ},
    \end{equation}
    and we take $q$ fixed, $N\to\infty$ (with the usual SYK scaling $\sigma\sim N
    ^{-(q-1)}$ understood).

    %------------------------------------------------------------

    In models in which the operators $V_{I}$ are (approximately) commuting at large
    system size---e.g., because most pairs $V_{I},V_{J}$ have disjoint support and
    hence commute while commutators require rare overlaps---the commutator corrections
    are suppressed. In the diffusive scaling regime
    \begin{equation}
        \tau \;:=\;\sigma t^{2}\binom{N}{q}\sim N t^{2}= O(1),
    \end{equation}
    the leading contribution $\Omega_{1}=O(1)$ survives while $\Omega_{n\ge2}$
    are parametrically smaller, so to leading order
    \begin{equation}
        M_{t}^{(k)}\;\simeq\; \exp\!\left[-\frac{\sigma t^{2}}{2}\sum_{I}\big(\widehat
        V_{I}^{(k)}\big)^{2}\right]. \label{eq:quenched_gaussian_regulator}
    \end{equation}

    %------------------------------------------------------------
    \paragraph{Identification of the effective generator.}
    Define
    \begin{equation}
        \widehat X_{I}\;:=\;\sqrt{\sigma}\,\widehat V_{I}^{(k)}=\sqrt{\sigma}\sum
        _{a=1}^{k}\!\left(V_{I,L}^{(a)}-\big(V_{I,R}^{(a)}\big)^{T}\right), \qquad
        \widehat G_{eff}^{(k)}\;:=\;\frac{1}{2}\sum_{I}\widehat X_{I}^{\,2}. \label{eq:hatGeff_quenched}
    \end{equation}
    Then~\eqref{eq:quenched_gaussian_regulator} becomes the Gaussian regulator
    \begin{equation}
        M_{t}^{(k)}\simeq e^{-\tau^2\hat{G}_{eff}^{(k)}}. \label{eq:Mt_exp_hatGeff}
    \end{equation}

    On the physical $k$-replica space define Hermitian operators
    \begin{equation}
        X_{I}\;:=\;\sqrt{\sigma}\sum_{a=1}^{k}V_{I}^{(a)}, \qquad \widehat X_{I}=
        (X_{I})_{L}- (X_{I})_{R}^{T}.
    \end{equation}
    Using
    \begin{equation}
        \big((X)_{L}-(X)_{R}^{T}\big)^{2}|Y\rangle\rangle = |X^{2}Y+Y X^{2}-2XYX\rangle
        \rangle,
    \end{equation}
    we obtain the corresponding superoperator $G_{eff}^{(k)}$ on $Y$ via
    \begin{equation}
        \widehat G_{eff}^{(k)}|Y\rangle\rangle = |G_{eff}^{(k)}(Y)\rangle\rangle,
    \end{equation}
    namely
    \begin{align}
        G_{eff}^{(k)}(Y) & =\frac{1}{2}\sum_{I}\Big(X_{I}^{2}Y+Y X_{I}^{2}-2X_{I}Y X_{I}\Big) \nonumber                                                 \\
                         & =\frac{1}{2}\sum_{I}[X_{I},[X_{I},Y]] \;=\;-\sum_{I}\Big(X_{I}Y X_{I}-\tfrac12\{X_{I}^{2},Y\}\Big). \label{eq:Geff_quenched}
    \end{align}
    Therefore, in the same large-$N$ regime in which~\eqref{eq:Mt_exp_hatGeff} holds,
    \begin{equation}
        \boxed{\mathcal{E}_t^{(k)} \;\simeq\; e^{-\tau^2\,G_{eff}^{(k)}}}
    \end{equation}

    \NL{Compute the thermal 2pt function, 4pt function and the multi-point OTOCs starting from the k-twirl map, at large N and various limits of $\beta J$.}

    \section{Gaussian regulator in large $N$ SYK}

    The SYK model is comprised of $N$ Majorana fermions $\chi_{i}$ with
    $i=1,\cdots , N$ and the Hamiltonian
    \begin{eqnarray}
        H=i^{q/2}\sum_{i_1<\cdots <i_q}J_{i_1<\cdots<i_q} (\chi_1\cdots \chi_q)\ .
    \end{eqnarray}
    In comparing to the notation of the previous section, $I$ is all ordered $q$-tuples
    $I = \{i_{1}< \dots < i_{q}\}$, and
    \begin{eqnarray}
        V_{i_1<\cdots <i_q}=\chi_{i_1}\chi_{i_2}\cdots \chi_{i_q}
    \end{eqnarray}
    is the product of $q$ Majorana fermions. It follows from $\{\chi_{i},\chi_{j}
    \}=2\delta_{ij}$ that
    \begin{eqnarray}
        (V_{i_1<\cdots <i_q})^2=i^q(\chi_{i_1}\cdots \chi_{i_q})^2=i^q(-1)^{q(q-1)/2}
    \end{eqnarray}
    which explains the origin of $i^{q/2}$ in (\ref{SYKHam}). In this note, we will
    restrict to the case where both $N$ and $q$ are even, and
    $(V_{i_1<\cdots <i_q})^{2}=1$.

    The time-independent noise is
    \begin{eqnarray}
        &&\overline{J_I}=0\nn\\
        &&\overline{J_I J_J}=\sigma \delta_{IJ}\nn\\
        &&\sigma=\frac{2^{q-1}q!}{q^{2}N^{q-1}}J\ .
    \end{eqnarray}
    Note that we are using the normalization in
    \cite{jian2021note,jian2023linear}. Note that there are $\binom{N}{q}$ terms
    in the Hamiltonian which in the limit of $q\ll N$ grows like $N^{q}/q!$. The
    factor $q!/N^{q-1}$ is designed to make sure that the fluctuations of the
    Hamiltonian are order $N$. For $q=2$ and $q=4$, we have, respectively $\sigma
    =J/N$ and $\sigma=12\,J/N^{3}$.

    \subsection{Jordan-Wigner transform}

    Assuming that $N$ is even, we pair Weyl fermions of each copy $\psi^{(a)}_{2j-1}$
    and $\psi^{(a)}_{2j}$ into $N/2$ complex fermions
    \begin{eqnarray}
        &&c_j^{(a)}\equiv \frac{1}{2}(\chi_{2j-1}^{(a)}+i \chi_{2j}^{(a)})\nn\\
        &&\{c_l,c_j^\dagger\}=\delta_{l j},\qquad \{c_l,c_j\}=\{c_l^\dagger,c_j^\dagger\}=0\ .
    \end{eqnarray}
    For a fixed $j$, the algebra generated by $c_{k},c_{k}^{\dagger}$ is
    isomorphic to the algebra of a qubit generated with
    $c_{j}=\ket{0}\bra{1}_{j}$ and we can identify $\chi_{2j-1}=X_{j}$ and
    $\chi_{2j}=Y_{j}$, where $X_{j}$ and $Y_{j}$ are Pauli matrices. The fermion
    number operator is
    \begin{eqnarray}
        N_j=c_j^\dagger c_j=\ket{1}\bra{1}_j=(1+Z_j)/2\ .
    \end{eqnarray}
    However, the problem is that $[X_{l},Y_{j}]=0$, whereas the Weyl fermions
    anticommute. To fix this, we define the operators
    \begin{eqnarray}
        &&\sigma^-_j\equiv \frac{1}{2}(X_j+i Y_j)=(-1)^{\sum_{l=1}^{j-1}c_l^\dagger c_l}c_j\nn\\
        &&Z_j=2c_j^\dagger c_j-1
    \end{eqnarray}
    by attaching to $c_{j}$ a string of parity operators
    \begin{eqnarray}
        (-1)^{N_j}=\ket{0}\bra{0}_j-\ket{1}\bra{1}_j\ .
    \end{eqnarray}
    It can now be checked explicitly that Paul matrices with different $j$
    commute. In terms of the Weyl fermions, we have the identification
    \begin{eqnarray}
        &&Z_j=-i \chi_{2j-1}\chi_{2j}=(-1)^{c_j^\dagger c_j+1}\nn\\
        &&X_j=\lb \prod_{l=1}^{j-1}(i \chi_{2l-1}\chi_{2l})\rb \chi_{2j-1}=(-1)^{\sum_{j=1}^{k-1}c_j^\dagger c_j}(c_k+c_k^\dagger)\nn\\
        &&Y_j=\lb \prod_{l=1}^{j-1}(i \chi_{2l-1}\chi_{2l})\rb \chi_{2j}=-i (-1)^{\sum_{l=1}^{j-1}c_l^\dagger c_l}(c_j-c_j^\dagger)
    \end{eqnarray}
    Equivalently, we can write
    \begin{eqnarray}
        &&c_j=(-1)^{j-1}(Z_1\cdots Z_{j-1})\sigma_j^-\nn\\
        &&\chi_{2j-1}=(-1)^{j-1}(Z_1\cdots Z_{j-1})X_j\nn\\
        &&\chi_{2j}=(-1)^{j-1}(Z_1\cdots Z_{j-1})Y_j\ .
    \end{eqnarray}

    Now, consider the following two purifications of a maximally mixed state of a
    single qubit (a single complex fermion) in a two-qubit (a pair of
    {\it commuting} complex fermions) system:
    \begin{eqnarray}
        &&\ket{e}=\frac{1}{\sqrt{2}}(\ket{00}+\ket{11})\nn\\
        &&\ket{f}=\frac{1}{\sqrt{2}}(\ket{01}+\ket{10})\nn\\
        &&c_{1L}=\sigma_{1L}^-,\qquad c_{2R}=\sigma_{1R}^-
    \end{eqnarray}
    which can be identified as the unique vector that satisfies the following
    relations
    \begin{eqnarray}
        &&(c_{1L}-c_{1R}^\dagger)\ket{e}=0=(c_{1L}^\dagger-c_{1R})\ket{e}\nn\\
        &&(c_{1L}-c_{1R})\ket{f}=0=(c_{1L}^\dagger-c_{1R}^\dagger)\ket{f}
    \end{eqnarray}
    Next, consider $N/2$ ($N$ even) pairs of Weyl fermions, and the maximally
    mixed vectors
    \begin{eqnarray}
        &&\ket{e}^{\otimes N/2},\qquad \ket{f}^{\otimes N/2}\nn\\
        &&c_{k,L/R}=(-1)^{k-1}Q_{k-1,L/R}\sigma_{k,L/R}^-,\nn\\
        &&Q_{k,L/R}\equiv (Z_{1L/R}\cdots Z_{k,L/R})\ .
    \end{eqnarray}
    Note that even though we are labeling the fermions using $L$ and $R$ they
    are commuting variables
    \begin{eqnarray}
        [c_{jL},c^\dagger_{l,R}]=0,\qquad \{c_{jL},c^\dagger_{lL}\}=\delta_{jl},\qquad
        \{c_{jR},c^\dagger_{lR}\}=\delta_{jl}
    \end{eqnarray}
    Since
    \begin{eqnarray}
        (Q_{j,L}-Q_{j,R})\ket{e}^{\otimes j}=0,\qquad (Q_{j,L}+(-1)^jQ_{j,R})\ket{f}^{\otimes k}=0
    \end{eqnarray}
    we find that for the following equations for all $l=1,\cdots, N/2$ for $N$
    even define the vectors $\ket{f}^{\otimes N/2}$ and $\ket{f}^{\otimes N/2}$:
    \begin{eqnarray}
        &&(c_{jL}-c_{jR}^\dagger)\ket{e}^{\otimes N/2}=0=(c_{jL}^\dagger-c_{jR})\ket{e}^{\otimes N/2}\nn\\
        &&(c_{jL}-c_{jR})\ket{f}^{\otimes N/2}=0=(c_{j L}^\dagger-c_{j R}^\dagger)\ket{f}^{\otimes N/2}\ .
    \end{eqnarray}
    A convenient way to relate the complex fermions to Majoranas is the following
    definitions\footnote{Note the difference compared to \cite{gu2017spread}. See
    their equation 2.12.}:
    \begin{eqnarray}
        c_{kL}=\frac{1}{2}(\chi_{2k-1,L}+i \chi_{2k,L}),\qquad c_{kR}=\frac{1}{2}(\chi_{2k,R}+i\chi_{2k-1,R})
    \end{eqnarray}
    This ensures that for all $j$ we have\footnote{Alternatively, we could use the
    convention in eq 2.12 of \cite{gu2017spread}
    \begin{eqnarray}
        c_{kL}=\frac{1}{2}(\chi_{2k-1,L}+i \chi_{2k,L}),\qquad c_{kR}=\frac{1}{2}(\chi_{2k,R}-i\chi_{2k-1,R})
    \end{eqnarray}
    that results in $(\chi_{j,L}+i\chi_{j,R})\ket{f}^{\otimes N/2}=0$.}
    \begin{eqnarray}
        \label{evector} (\chi_{j,L}+i\chi_{j,R})\ket{e}^{\otimes N/2}=0\ .
    \end{eqnarray}
    This is a very convenient choice, because the terms in the Hamiltonian involve
    all $\chi_{j}$, and without the need to distinguish the even and odd labels,
    for even $q$ we can write
    \begin{eqnarray}
        &&(H\otimes 1)\ket{e}^{\otimes N/2}=i^{q/2}\sum_{i_1<\cdots i_q}J_{i_1\cdots i_q}\chi_{i_1L}\cdots
        \chi_{i_qL}\ket{e}^{\otimes N/2}\nn\\
        &&=i^{q/2}\sum_{i_1<\cdots i_q}J_{i_1\cdots i_q}\chi_{i_1R}\cdots \chi_{i_qR}\ket{e}^{\otimes N/2}=(1\otimes
        H)\ket{e}^{\otimes N/2}\ .
    \end{eqnarray}
    For a general operator $\mO_{L}$, the mirror operator in the vector
    $\ket{e}^{\otimes N/2}$ is $\mO^{T}_{R}$
    \begin{eqnarray}
        (\mO\otimes 1_R)\ket{e}^{\otimes k}=(1_L\otimes \mathcal{O}^T) \ket{e}^{\otimes k}\ .
    \end{eqnarray}
    This means in the qubit basis, $H$ is symmetric in the qubit basis. The
    advantage of the qubit basis is that it is totally commuting, which is convenient
    for applying quantum circuit models. However, the Hamiltonian expressed in terms
    of the qubit degrees of freedom does not take a simple form. To address this,
    in the next subsection, following \cite{gu2017spread}, we use totally anti-commuting
    variables, which have the advantage that the Hamiltonian $\hat{H}$ takes a
    simple form in these variables. Bilinears of these fermions are commuting variables
    and well-suited for quantum circuit models.

    \subsection{Totally anti-commuting variables}

    Since we are interested in $k$-twirl map $(U\otimes U^{*})^{\otimes k}$, consider
    $2k$ copies of SYK fermions that we label by $\chi^{(a)}_{L j}$ and $\chi^{(a)}
    _{R j}$ with $j=1,\cdots, N$ with $N$ even, $a=1,\cdots, k$, and $L$ and $R$
    denoting where they are evolved by $U$ or $U^{\dagger}$.

    We define the total parity operator, which counts the total number of
    complex fermions for a fixed $(a)$:
    \begin{eqnarray}
        Q^{(a)}_{L/R}=(-1)^{\sum_{j=1}^{N/2}n_j}=\prod_{j=1}^{N/2}Z_{j (L/R)}=(-i)^{N/2}\prod_{l=1}^{N}\chi_{l (L/R)}^{(a)}
    \end{eqnarray}
    For even $N$ it satisfies the following relations
    \begin{eqnarray}
        &&\{Q^{(a)}_{L/R},\chi^{(a)}_{l,L/R}\}=0,\qquad a\neq b: [Q^{(a)}_{L/R},\chi^{(b)}_{l,L/R}]=0,\nn\\
        &&[Q^{(a)}_L,\chi^{(b)}_{l,R}]=0,\qquad (Q^{(a)}_{L/R})^2=1^{(a)}_{L/R}\ .
    \end{eqnarray}
    Then, we define new Majorana fermions \footnote{Note the asymmetry in the
    definition between $L$ and $R$; i.e. the extra factor of $i$.}
    \begin{eqnarray}
        &&\psi_{l,L}^{(a)}=i \Sigma^{(a)}\chi_{l,L}^{(a)}, \qquad \psi_{l,R}^{(a)}=\Sigma^{(a)}\chi_{l,R}^{(a)},\nn\\
        &&\Sigma^{(a)}=\lb\prod_{b=1}^{a-1} Q_L^{(h)} Q_R^{(b)}\rb Q^{(a)}_L
    \end{eqnarray}
    that anti-commute for all labels
    \begin{eqnarray}
        \{\psi^{(a)}_{l, \alpha},\psi^{(b)}_{l',\beta}\}=2\delta_{ll'}\delta_{ab}\delta_{\alpha\beta}
    \end{eqnarray}
    with $\alpha,\beta=L,R$. The Majorana condition follows from the (anti)-commutation
    relations of $Q^{(a)}_{L/R}$ with $\chi$ Majorana fermions:
    \begin{eqnarray}
        (\psi^{(a)}_{l,L})^\dagger=\psi^{(a)}_{l,L},\qquad (\psi^{(a)}_{l,R})^\dagger=\psi^{(a)}_{l,R}\ .
    \end{eqnarray}
    It follows from (\ref{evector}) that $\psi_{l,L}^{(a)}$ are
    $\psi^{(a)}_{l,R}$ are mirror operators in this vector:
    \begin{eqnarray}
        (\psi^{(a)}_{l,L}-\psi^{(a)}_{l,R})\ket{e}_a^{\otimes N/2}=0\ .
    \end{eqnarray}
    A nice feature of the new Majorana $\psi$ is that for all $i_{1}<i_{2}<\cdots
    <i_{q}$ with $q$ even we have
    \begin{eqnarray}
        \chi^{(a)}_{i_1,L}\cdots \chi^{(a)}_{i_q,L}=\psi^{(a)}_{i_1,L}\cdots \psi^{(a)}_{i_q,L}\ .
    \end{eqnarray}
    Therefore, we have
    \begin{eqnarray}
        \hat{V}^{(a)}_{i_1\cdots i_q}=i^{q/2}\psi^{(a)}_{i_1,\alpha}\cdots \psi^{(a)}_{i_q,\alpha}
    \end{eqnarray}
    It is straightforward to check explicitly that each term $H^{(a)}_{I}$ in
    the Hamiltoian saistifes the mirror equation:
    \begin{eqnarray}
        \psi^{(a)}_{i_1L}\cdots \psi^{(a)}_{i_qL}\ket{e}_a^{\otimes N/2}= \psi^{(a)}_{i_1R}\cdots
        \psi^{(a)}_{i_qR}\ket{e}_a^{\otimes N/2}\ .
    \end{eqnarray}
    This means that in the case of SYK, we do not need to worry about transposes,
    because every term in the Hamiltonian is symmetric in the basis of $\ket{e}^{\otimes
    N/2}$. Therefore, we can define
    \begin{eqnarray}
        \label{VSYK} \hat{V}^{(a)}_{i_1<\cdots <i_q}= i^{q/2}(\psi^{(a)}_{i_1L}\cdots
        \psi^{(a)}_{i_qL}- \psi^{(a)}_{i_1R}\cdots \psi^{(a)}_{i_qR})\ .
    \end{eqnarray}

    %\NL{How about the $g$-averaged higher point functions?}

    % A simple trick to compute higher point functions in the (unnormalized) infinite temperature thermofield double $\ket{1}$ is to compute
    % \begin{eqnarray}
    %   &&\mathcal{E}_g(U^{\otimes k}\otimes (U^*)^{\otimes k})(\mathbb{Y}\otimes 1^{\otimes k})\ket{1}^{\otimes k}\nn\\
    %   &&=\mathcal{E}_g\lb e^{-i g \hat{\mathbb{H}}}\rb(\mathbb{Y}\otimes 1^{\otimes k})\ket{1}^{\otimes k}\nn\\
    %   &&\hat{\mathbb{H}}=\sum_I g_I \sum_a \hat{V}_I^{(a)}\nn\\
    %   &&\hat{\mathbb{X}}_I=\sqrt{\sigma}\sum_a \hat{V}_I^{(a)}\ .
    %   %\rb\nn\\
    %   %&&=\tr((\mO_1\otimes \cdots \otimes\mO_k)W_{\pi_{cyc}})
    % \end{eqnarray}
    % where in our notation $\mathbb{Y}$, $\mathbb{X}_I$ and $\mathbb{H}_I$ acts on $\mathcal{H}^{\otimes k}$, and their hatted versions $\hat{\mathbb{H}}$, $\hat{\mathbb{X}}$ and $\hat{\mathbb{H}}$ act on $\mathcal{H}^{2k}$. Performing the Gaussian integral as before, we find
    % \begin{eqnarray}
    %   &&\mathcal{E}_g(U^{\otimes k}\otimes (U^*)^{\otimes k})=e^{-\hat{\mathbb{G}}_{eff}t^2}\nn \\
    %   &&\hat{\mathbb{G}}_{eff}=\frac{1}{2}\sum_I \hat{\mathbb{X}}_I^2\ .
    % \end{eqnarray}
    % This GNS operator can be viewed as a superoperator acting on operators in $\mathcal{H}^{\otimes k}$:
    % \begin{eqnarray}
    %   &&e^{-\hat{\mathbb{G}}_{eff}^{(k)}t^2}(\mathbb{Y}\otimes 1^{\otimes k})\ket{1}^{\otimes k}=(\Phi_t^{(k)}(\mathbb{Y})\otimes 1^{\otimes k})\ket{1}^{\otimes k}\nn\\
    %   &&\hat{\mathbb{G}}_{eff}^{(k)}(\mathbb{Y}\otimes 1)\ket{1}^{\otimes k}=(\mathcal{G}_{eff}^{(k)}(\mathbb{Y})\otimes 1^{\otimes k})\ket{1}^{\otimes k}\nn\\
    %   &&\mathcal{G}_{eff}^{(k)}(\mathbb{Y})=\lb \mathbb{X}_I \mathbb{Y}\mathbb{X}_I^\dagger-\frac{1}{2}(\mathbb{X}_I^2\mathbb{Y}+\mathbb{Y}(\mathbb{X}_I^\dagger)^2\rb\nn\\
    %     &&=\lb \mathbb{X}_I \mathbb{Y}\mathbb{X}_I-\frac{1}{2}\{\mathbb{X}_I^2,\mathbb{Y}\}\rb
    %   \end{eqnarray}
    %   where we have used the fact that each $\mathbb{X}_I$ term is self-adjoint. The result
    %   resembles a Linbladian with self-adjoint jump operators $\mathbb{X}_I$ and decay rates $\sigma$, except that it sits next to $t^2$ instead of $t$! Of course, this is NOT a semi-group, and the $k$-fold quantum channel $\Phi_t$ can be written as
    %   \begin{eqnarray}
    %     \Phi_{t}^{(k)}(\mathbb{Y})=\sum_{n=0}^\infty \frac{(-t^2)^n}{n!}(\mathcal{G}_{eff}^{(k)})^n(\mathbb{Y})\ .
    %   \end{eqnarray}
    % where $W_{\pi_{cyc}}$ is the cyclic permutation $i\to (i+1)\mod k$ on $\mathcal{H}^{\otimes k}$.

    % Note that, in general, $[V_I,V_J]\neq 0$, therefore $[G_{eff},H(g)]\neq 0$ and $[G_{eff},\rho_\beta]=0$.

    % \NL{A convenient limit here is to take the infinite temperature limit $\beta\to 0$, keeping $g$ and $\sigma$ finite. This amounts to sending $\tilde{g}\to 0$, and $\tilde{\sigma}\to 0$}

    %------------------------------------------------------------

    \subsection{$k$-twirl map in large $N$ SYK}

    The effective generator we found is
    \begin{eqnarray}
        \allowdisplaybreaks &&\mathbb{G}_{eff}^{(k)}=(-1)^{q/2}\sum_{i_1<\cdots <i_q}(\hat{\mathbb{X}}_{i_1<\cdots <i_q})^2\nn\\
        &&=(-1)^{q/2}\sum_{i_1<\cdots <i_q}\sum_{a,b=1}^k(\psi_{i_1L}^{(a)}\cdots
        \psi_{i_qL}^{(a)}-\psi_{i_1R}^{(a)}\cdots \psi_{i_qR}^{(a)})(\psi_{i_1L}^{(b)}\cdots
        \psi_{i_qL}^{(b)}-\psi_{i_1R}^{(b)}\cdots \psi_{i_qR}^{(b)})\nn\\
        &&=(-1)^{q/2}\sum_{i_1<\cdots <i_q}\sum_{a,b=1}^k\sum_{\alpha,\beta=L,R}(-1)^{q(q-1)/2}\lb
        \psi^{(ab)}_{i_1;\alpha\beta}\cdots \psi^{(ab)}_{i_q;\alpha\beta}\rb\nn\\
        &&=(-i)^q\sum_{i_1<\cdots <i_q}\sum_{a,b=1}^k\sum_{\alpha,\beta=L,R}\lb
        \psi^{(ab)}_{i_1;\alpha\beta}\cdots \psi^{(ab)}_{i_q;\alpha\beta}\rb\label{Geffk}
    \end{eqnarray}
    where we have defined the local bilinears
    \begin{eqnarray}
        \label{bilinearpsi} &&M^{(ab)}_{l;\alpha\beta}\equiv i\psi^{(a)}_{l \alpha}\psi^{(b)}_{l\beta}\ .
    \end{eqnarray}
    used that for even $q$, the factor $(-1)^{q/2}(-1)^{q(q-1)/2}=1$.
    % , and that for commutative $A_i$ we have
    % \begin{equation}
    % \sum_{i_1<\cdots <i_q}A_{i_1}\cdots A_{i_q}=\frac{1}{q!}\sum_{\text{distinct} \: i_1\cdots i_q}A_{i_1}\cdots A_{i_q}
    % \end{equation}
    Note that $M^{(aa)}_{l;\alpha\alpha}=1$, and

    The bilinears at different sites $j\neq j'$ commute:
    \begin{eqnarray}
        [M^{(ab)}_{j;\alpha\beta},M^{(a'b')}_{j';\alpha'\beta'}]=0
    \end{eqnarray}
    Next, note that since $\psi_{j;LR}^{(ab)}=-\psi_{j;RL}^{(ab)}$ using the
    fact that $q$ is even, we have
    \begin{eqnarray}
        &&\mathbb{G}_{eff}^{(k)}= \sum_{a,b=1}^k\lb e_{q;LL}^{(ab)}+e_{q;RR}^{(ab)}-2e_{q;LR}^{(ab)}\rb\\
        &&e_{q;\alpha\beta}^{(ab)}=\sum_{i_1<\cdots <i_q}\prod_{m=1}^q\psi^{(ab)}_{i_m;\alpha\beta}
    \end{eqnarray}

    In the $k$-fold (twirl) construction it is often convenient to bundle the
    replica label $a=1,\dots,k$ and the side label $\alpha\in\{L,R\}$ into a
    single ``flavor'' index\footnote{Equivalently, fix any bijection between $(a,
    \alpha)$ and $\{1 , \dots,2k\}$; an ordering convention then allows
    shorthand sums such as $\sum_{A<B}$ to run over unordered pairs once.}
    \begin{eqnarray}
        &&A \equiv (a,\alpha)\,,\qquad A=1,\dots,2k\nn\\
        &&s_A=
        \begin{cases}
             & 1\qquad A=L    \\
             & (-1)\qquad A=R
        \end{cases}\ .
    \end{eqnarray}
    Our Majorana fermions $\psi_{i}^{A}$ labeled by sites $i=1,\dots,N$ satisfy
    \begin{equation}
        \{\psi_{i}^{A},\psi_{j}^{B}\}=2\,\delta_{ij}\,\delta^{AB}\,, \qquad (\psi
        _{i}^{A})^{\dagger}=\psi_{i}^{A}\,. \label{eq:majorana_CAR}
    \end{equation}
    Motivated by (\ref{bilinearpsi}) we define the antisymmetric bilinears (sitewise
    generators)
    \begin{equation}
        M_{i}^{AB}\;\equiv\;\frac{i}{2}\,[\psi_{i}^{A},\psi_{i}^{B}]\,,\qquad M_{i}
        ^{AB}=-M_{i}^{BA}\ . \label{eq:def_M}
    \end{equation}
    Using \eqref{eq:majorana_CAR}, for $A\neq B$ we have $[\psi_{i}^{A},\psi_{i}^{B}
    ]=2\psi_{i}^{A}\psi_{i}^{B}$, so the operator $M^{AB}_{i}$ are self-adjoint:
    \begin{equation}
        \boxed{ M_i^{AB}= i\,\psi_i^A\psi_i^B \qquad (A\neq B)\,\qquad M_i^{AA}=0}
        \ . \label{eq:M_as_product}
    \end{equation}
    For $i\neq j$, the Majoranas anticommute, $\psi_{i}^{A}\psi_{j}^{B}=-\psi_{j}
    ^{B}\psi_{i}^{A}$. Since each $M_{i}^{AB}$ is an \emph{even} (quadratic) operator
    at site $i$, bilinears at distinct sites commute:
    \begin{equation}
        [M_{i}^{AB},M_{j}^{CD}] = 0\,,\qquad i\neq j\,. \label{eq:M_commute_different_sites}
    \end{equation}
    A direct proof uses \eqref{eq:M_as_product} (for $A\neq B$ and $C\neq D$):
    \begin{align}
        M_{i}^{AB}M_{j}^{CD} & = (i\psi_{i}^{A}\psi_{i}^{B})(i\psi_{j}^{C}\psi_{j}^{D}) = -\,\psi_{i}^{A}\psi_{i}^{B}\psi_{j}^{C}\psi_{j}^{D}\nonumber \\
                             & = -\,\psi_{j}^{C}\psi_{j}^{D}\psi_{i}^{A}\psi_{i}^{B}= M_{j}^{CD}M_{i}^{AB}\,,
    \end{align}
    where moving $\psi_{j}^{C}\psi_{j}^{D}$ past $\psi_{i}^{A}\psi_{i}^{B}$ produces
    $(-1)^{2\cdot 2}=+1$. Since we also know
    \begin{eqnarray}
        (M^{AB}_i)^2=\delta_{AB}\,
    \end{eqnarray}
    for $A\neq B$, they are self-adjoint unitary operators, therefore their
    spectrum is either $+1$ or $-1$.

    For $i\neq j$ one also has
    \begin{equation}
        [M_{i}^{AB},\psi_{j}^{C}]=0\,,\qquad i\neq j\,, \label{eq:M_commute_fermion_other_site}
    \end{equation}
    since moving $\psi_{j}^{C}$ past the quadratic operator at site $i$ produces
    $(-1)^{2}=+1$. Whereas for the same site ($i=j$), we have
    \begin{align}
        [\psi_{i}^{A}\psi_{i}^{B},\psi_{i}^{C}] & = \psi_{i}^{A}\{\psi_{i}^{B},\psi_{i}^{C}\}-\{\psi_{i}^{A},\psi_{i}^{C}\}\psi_{i}^{B}\nonumber                         \\
                                                & = \psi_{i}^{A}(2\delta^{BC})-(2\delta^{AC})\psi_{i}^{B}= 2\big(\delta^{BC}\psi_{i}^{A}-\delta^{AC}\psi_{i}^{B}\big)\,.
    \end{align}
    Multiplying by $i$ gives the compact result
    \begin{equation}
        \boxed{ [M_i^{AB},\psi_j^C] \;=\; 2i\,\delta_{ij}\,\big(\delta^{BC}\psi_i^A-\delta^{AC}\psi_i^B\big)\,. }
        \label{eq:M_comm_psi}
    \end{equation}
    This explicitly shows that $M_{i}^{AB}$ generates orthogonal rotations of
    the flavor index on the Majoranas at each site.
    %  Since $A,B=1,\cdots, 2k$, we have the Lie-algebra of $\mathfrak{so}(2k)$, corresponding to linear Bogoliubov transformations of the fermions $\psi_i^A$.

    % \subsection{Bilinear--bilinear commutators (closure)}
    % \label{subsec:M_comm_M}
    At the same site, the $M_{i}^{AB}$ close into the $\mathfrak{so}(2k)$ Lie algebra.
    Starting from
    \begin{equation}
        [\psi_{i}^{A}\psi_{i}^{B},\psi_{i}^{C}\psi_{i}^{D}] = [\psi_{i}^{A}\psi_{i}
        ^{B},\psi_{i}^{C}]\psi_{i}^{D}+ \psi_{i}^{C}[\psi_{i}^{A}\psi_{i}^{B},\psi
        _{i}^{D}]\,,
    \end{equation}
    and inserting the previously derived single-fermion commutator results in
    \begin{equation}
        \boxed{ [M_i^{AB},M_j^{CD}] = 2i\,\delta_{ij}\Big( \delta^{BC}M_i^{AD} -\delta^{AC}M_i^{BD} -\delta^{BD}M_i^{AC} +\delta^{AD}M_i^{BC} \Big)\,. }
        \label{eq:so2k_commutator}
    \end{equation}
    Together with \eqref{eq:M_commute_different_sites}, Eq.~\eqref{eq:so2k_commutator}
    shows that each site carries an independent copy of $\mathfrak{so}(2k)$
    corresponding to the linear Bogoliubov transformations of fermions
    $\psi_{i}^{A}$, while different sites commute.

    % \paragraph{Remark (ordering and $\sum_{A<B}$).}
    % Because $M_i^{AB}$ is antisymmetric in $(A,B)$ and $M_i^{AA}=0$, it is often convenient to sum
    % only over unordered pairs once:
    % \begin{equation}
    % \sum_{A<B}(\cdots)\equiv \frac{1}{2}\sum_{A\neq B}(\cdots)\,,
    % \end{equation}
    % where the inequality $A<B$ refers to any fixed total ordering of the $2k$ flavors.

    Now, we can write the expression of the effective Hamiltonian in (\ref{Geffk})
    in terms of these bilinears:
    \begin{eqnarray}
        \label{GeffkM} \mathbb{G}_{eff}^{(k)}&&=i^q\sum_{i_1<\cdots <i_q}\lb \sum_A
        1+s_As_B\sum_{AB}M_{i_1}^{AB}\cdots M_{i_q}^{AB}\rb\nn\\
        &&=i^q\lb \sum_A 2k{N\choose q}+\sum_{i_1<\cdots <i_q}s_As_B\sum_{AB}M_{i_1}^{AB}\cdots
        M_{i_q}^{AB}\rb
    \end{eqnarray}
    Note that since $M_{i}^{aL,bR}=-M_{i}^{bL,aR}$ and $q$ is even, we can write
    \begin{eqnarray}
        \mathbb{G}_{eff}^{(k)}=i^q\lb \sum_A 2k{N\choose q}+\sum_{i_1<\cdots <i_q}\sum_{ab}\lb
        M_{i_1,LL}^{ab}\cdots M_{i_q,LL}^{ab}+M_{i_1,RR}^{ab}\cdots M_{i_q,RR}^{ab}-2M_{i_1,LR}^{ab}\cdots
        M_{i_q,LR}^{ab}\rb\rb
    \end{eqnarray}
    We can also define new generators
    \begin{eqnarray}
        \tilde{M}^{AB}_i\equiv s_As_BM^{AB}_i
    \end{eqnarray}
    which also generate $\mathfrak{so}(2k)$ algebra because this redefition is congugation
    by an $O(2k)$ matrix. To interpret the generator (\ref{}) it is convenient
    to write it in terms of
    \begin{eqnarray}
        \label{eABq} &&\tilde{M}^{AB}=\sum_{i=1}^N \tilde{M}^{AB}_i\nn\\
        &&e^{AB}_q=\sum_{i_1<\cdots <i_q}\tilde{M}^{AB}_{i_1}\cdots \tilde{M}^{AB}_{i_q}=\sum_{r=0}^{\floor{q/2}}c_{q,r}(N)(\tilde{M}^{AB})^{q-2r}
    \end{eqnarray}
    using a generating function $\sum_{q=0}^{N}z^{q}e^{AB}_{q}$. Here, the
    coefficients $c_{q,r}(N)$ are called the Krawtchouk polynomials starting
    with
    \begin{eqnarray}
        &&e_2^{AB}=\frac{(\tilde{M}^{AB})^{2}-N}{2},\nn\\
        &&e_4^{AB}=\frac{(\tilde{M}^{AB})^{4}-2(3N-4)(\tilde{M}^{AB})^{2}+3N(N-2)}{2}\ .
    \end{eqnarray}

    Note that the operators $\tilde{M}_{i}^{AB}$ are self-adjoint and square to identity,
    their spectrum must be $(-1,+1)$. Then, since $\tilde{M}^{AB}$ is a sum of
    $N$ commuting $\tilde{M}_{i}^{AB}$, its spectrum follows the binomial distribution
    \begin{eqnarray}
        S=\{-N,-N+2,\cdots , N-2,N\}\ .
    \end{eqnarray}
    Plugging this in the expresison for $e^{AB}_{q}$ (\ref{eABq}) we find that it
    has the spectrum
    \begin{eqnarray}
        e^{AB}_q(S)=\sum_{r=0}^{\floor{q/2}}c_{q,r}(N)S^{q-2r}
    \end{eqnarray}

    It is convenient to strip off factorials and define monic polynomials
    $p_{q}( S)\equiv q!\,e_{q}^{AB}(S)$. Using Newton identities together with $p
    _{2m}=\sum_{i}(\tilde M_{i}^{AB})^{2m}=N$ and $p_{2m+1}=\sum_{i}(\tilde M_{i}
    ^{AB})^{2m+1}=S^{AB}$, one finds the three-term recursion
    \begin{eqnarray}
        \label{eq:krawtchouk_recursion} p_{q+1}(S)=S\,p_q(S)-q\,(N-q+1)\,p_{q-1}(S),\qquad
        p_0(S)=1,\ \ p_1(S)=S\ .
    \end{eqnarray}
    This is the standard Krawtchouk (Jacobi) recursion; equivalently, $p_{q}(S)$
    is the characteristic polynomial of the $q\times q$ tridiagonal matrix with
    diagonal entries $S$ and off-diagonal entries $\sqrt{m(N-m+1)}$ ($m=1,\dots,q
    -1$). Expanding $p_{q}(S)=\sum_{r=0}^{\lfloor q/2\rfloor}\tilde c_{q,r}(N)\,S
    ^{q-2r}$ and comparing with \eqref{eq:eq_expand_monomials} gives
    $c_{q,r}(N)=\tilde c_{q,r}(N)/q!$ and the coefficient recursion
    \begin{eqnarray}
        \label{eq:cqr_recursion} (q+1)\,c_{q+1,r}(N)=c_{q,r}(N)-(N-q+1)\,c_{q-1,r-1}(N),
    \end{eqnarray}
    with $c_{0,0}=1$, $c_{1,0}=1$, and $c_{q,r}=0$ for $r<0$ or
    $r>\lfloor q/2\rfloor$. A fully explicit closed form (weighted matchings on
    a path) is
    \begin{eqnarray}
        \label{eq:cqr_matching} c_{q,r}(N)=\frac{(-1)^{r}}{q!} \!\!\!\sum_{\substack{1\le m_1<\cdots<m_r\le q-1\\ m_{j+1}\ge m_j+2}}
        \ \prod_{j=1}^{r} \Big(m_j\,(N-m_j+1)\Big)\ .
    \end{eqnarray}
    In particular,
    \begin{eqnarray}
        c_{q,0}(N)=\frac{1}{q!},\qquad c_{q,1}(N)=-\frac{q(q-1)(3N-2q+4)}{6\,q!},
    \end{eqnarray}
    and for $q=2n$ and $N$ even the constant term is
    \begin{eqnarray}
        c_{2n,n}(N)=(-1)^n\binom{N/2}{n}\ .
    \end{eqnarray}

    \subsection{Properties of the generator $G_{eff}^{(k)}$}
    \subsubsection{Large-$N$ limits of $e_{q}^{AB}$}
    Fix a pair $A<B$ and denote
    \begin{eqnarray}
        S^{AB}\equiv \sum_{i=1}^N \tilde M^{AB}_i,\qquad m^{AB}\equiv
        \frac{S^{AB}}{N},\qquad e_q^{AB}\equiv \sum_{i_1<\cdots<i_q}\tilde M^{AB}_{i_1}\cdots
        \tilde M^{AB}_{i_q}\ .
    \end{eqnarray}
    Throughout we assume $N$ and $q$ are even. Since for fixed $(A,B)$ we have
    $[ \tilde M^{AB}_{i},\tilde M^{AB}_{j}]=0$ for $i\neq j$ and
    $(\tilde M^{AB}_{i})^{2}=\mathbbm 1$, the family $\{\tilde M^{AB}_{i}\}_{i=1}
    ^{N}$ is simultaneously diagonalizable with eigenvalues $\pm 1$, and
    $e_{q}^{AB}$ is a polynomial in the single collective variable $S^{AB}$.

    \paragraph{1) $N\to\infty$ with $q$ fixed.}
    In this regime it is natural to normalize by the number of $q$-subsets. On
    each joint eigenspace (so $m^{AB}\in[-1,1]$ is a number),
    \begin{eqnarray}
        \label{eq:eq_largeN_fixedq} \frac{e_{q}^{AB}}{\binom{N}{q}} = (m^{AB})^q-\frac{q(q-1)}{2N}\big(1-(m^{AB})^2\big)(m^{AB})^{q-2}
        +O\!\left(\frac{1}{N^{2}}\right)\ ,
    \end{eqnarray}
    and in particular $e_{q}^{AB}=\binom{N}{q}(m^{AB})^{q}+O(N^{q-1})$.

    \paragraph{2) $N,q\to\infty$ with $q^{2}/N$ fixed.}
    Let $q\sim \sqrt{N}$ so that $q^{2}/N=\lambda$ is held fixed. Using the generating
    function $\sum_{q\ge 0}z^{q}e_{q}^{AB}=(1+z)^{(N+S^{AB})/2}(1-z)^{(N-S^{AB})/2}$
    and the fact that the coefficient saddle sits at $z\sim q/N\sim N^{-1/2}$,
    we may expand the logarithms to quadratic order,
    \begin{eqnarray}
        (1+z)^{\frac{N+S^{AB}}{2}}(1-z)^{\frac{N-S^{AB}}{2}} =\exp\!\left(S^{AB}z-\frac{N}{2}z^2+o(1)\right),
    \end{eqnarray}
    which implies the Hermite scaling form
    \begin{eqnarray}
        \label{eq:eq_hermite_scaling} e_q^{AB} =\frac{N^{q/2}}{q!}\, \mathrm{He}_q\!\left(\frac{S^{AB}}{\sqrt{N}}\right)\,\big(1+o(1)\big)\ ,
    \end{eqnarray}
    where $\mathrm{He}_{q}$ are the probabilists' Hermite polynomials defined by
    \begin{eqnarray}
        e^{xt-\frac{t^{2}}{2}}=\sum_{q=0}^\infty \mathrm{He}_q(x)\,\frac{t^{q}}{q!}\ .
    \end{eqnarray}

    Then Eq.~\ref{GeffkM} can be written as
    \begin{eqnarray}
        \label{eq:Geff_poly_form} G^{(k)}_{eff} = i^{q}\Bigg(2k\binom{N}{q}+\sum_{A<B}
        e_q^{AB}\Bigg), \qquad e_q^{AB}=\sum_{i_1<\cdots<i_q}\tilde M^{AB}_{i_1}\cdots
        \tilde M^{AB}_{i_q}\ .
    \end{eqnarray}
    For each fixed $(A,B)$, the operators $\{\tilde M^{AB}_{i}\}_{i=1}^{N}$ are commuting
    involutions, hence $e_{q}^{AB}$ depends only on the single collective
    variable $\tilde M^{AB}$ and for $q$ even admits the expansion
    \begin{eqnarray}
        \label{eq:eq_poly_expand} e_q^{AB}(\tilde M^{AB})=\sum_{r=0}^{q/2} c_{q,r}(N)\,(\tilde
        M^{AB})^{q-2r}\ .
    \end{eqnarray}
    Substituting \eqref{eq:eq_poly_expand} into \eqref{eq:Geff_poly_form} yields
    \begin{eqnarray}
        \label{eq:Geff_sum_over_AB} G^{(k)}_{eff} = i^{q}\Bigg(2k\binom{N}{q} +\sum_{A<B}\sum_{r=0}^{q/2}
        c_{q,r}(N)\,(\tilde M^{AB})^{q-2r}\Bigg)\ .
    \end{eqnarray}

    \paragraph{Algebraic structure and limits.}
    At each site $i$, the bilinears $M^{AB}_{i}$ generate a copy of $\mathfrak{so}
    (2k)$, while different sites commute; therefore the global sums $\tilde M^{AB}
    =\sum_{i}\tilde M^{AB}_{i}$ generate the diagonal $\mathfrak{so}(2k)$ action
    and Eq.~\eqref{eq:Geff_sum_over_AB} exhibits $G^{(k)}_{eff}$ as an element of
    $U(\mathfrak{so}(2k))$ built as a sum over directions $(A,B)$ of the same degree-$q$
    polynomial in $\tilde M^{AB}$. For $N\to\infty$ with $q$ fixed, $e_{q}^{AB}=\binom
    {N}{q}(\tilde M^{AB}/N)^{q}+O(N^{q-1})$; in the scaling limit $N,q\to\infty$
    with $q^{2}/N$ fixed, $e_{q}^{AB}\sim \frac{N^{q/2}}{q!}\mathrm{He}_{q}(\tilde
    M^{AB}/\sqrt{N})$.

    Aa a superoperator the $k$-twirl map is a Lindbladian with self-adjoint jump
    operators $X_{I}$ (with $I\equiv i_{1}<\cdots<i_{q}$):
    \begin{eqnarray}
        \mathcal{G}^{(k)}_{eff}(Y) &&=\sum_I\!\left(\mathbb{X}_I\,Y\,\mathbb{X}_I-\frac{12}{\{}\mathbb{X}_I^2,Y\}\right)
        =-\frac{12}{\sum}_I [\mathbb{X}_I,[\mathbb{X}_I,Y]]\,,\nn\\
        &&=\sum_{I} \sum_{a,a'=1}^k\!\left(V_I^{(a)}\,Y\,V_I^{(a')}-\frac{12}{}V_I^{(a)}V_I^{(a')},Y\}\right),
    \end{eqnarray}
    where $V_{I}^{(a)}$ denotes the operator $V_{I}$ acting on replica $a$ (and as
    the identity on the others).

    \section{Semi-groups from time-depenent (Brownian) dynamics}

    To simplify our notation, from now on, we use dimensionless parameters and denote
    them by $\tilde{t}\to t$, $\tilde{\sigma}\to \sigma$, $\tilde{g}_{I}\to g_{I}$,
    and $\tilde{H}\to H$. We repeat the example above, but this time, splitting the
    total evolution time $T$ into $T/\epsilon$ intervals $(t_{i},t_{i}+\epsilon)$,
    and making the coupling to depend on the interval $g(t_{i})$:
    \begin{eqnarray}
        H_g(t_i)=\sum_I g_I(t_i)V_I\ .
    \end{eqnarray}
    We pick the coupling $g(t_{i})$ during this time interval by taking i.i.d.
    samples from a Gaussian distribution so that
    \begin{eqnarray}
        &&\mathcal{E}_g(g_I(t_i))=0\nn\\
        &&\mathcal{E}_g(g_I(t_i)g_J(t_j))=\frac{1}{\ep}\delta_{ij}\sigma
    \end{eqnarray}
    After one time interval $\ep$, we have have
    \begin{eqnarray}
        &&\mathcal{E}_g(dU^{\otimes k}\otimes (dU^*)^{\otimes k})=\mathcal{E}_g\lb
        1^{\otimes 2k}-i\ep \hat{\mathbb{H}}_g-\frac{\ep^{2}}{2}\hat{\mathbb{H}}_g^2+\cdots\rb
        \nn\\
        &&dU=U(t_i,t_i+\ep),\qquad \hat{\mathbb{H}}_g=\sum_{a=1}^k \hat{H}_g^{(a)},\qquad\hat{H}_g^{(a)}=H_{g,L}^{(a)}-(H_{g,R}^{(a)})^T\nn\ .
    \end{eqnarray}
    where the index $a=1,\cdots , k$ keeps track of the $k$ copies in the twirl map.
    Note that since different samples of $g_{I}(t_{i})$ result in operators that
    do not commute, we cannot pick a vector $\ket{1}$ that simulatenously diagonalizes
    them all. Therefore, for any fixed basis, we need to write
    $(H_{g,R}^{(a)})^{T}$ as the mirror operator of $H_{g,L}^{(a)}$.

    Since
    \begin{eqnarray}
        &&\mathcal{E}_g(\hat{\mathbb{H}_g})=\sum_{a I}\mathcal{E}(g_I)\hat{V}_I^{(a)}=0\nn\\
        &&\mathcal{E}_g(\hat{\mathbb{H}}_g^2)=\sum_{a,a',I,J}\mathcal{E}_g(g_Ig_J)\hat{V}^{(a)}_I
        \hat{V}_J^{(a')}=\frac{\sigma}{\ep}\sum_{a,a',I} \hat{V}_I^{(a)}\hat{V}_I^{(a')}\ .
    \end{eqnarray}
    We define
    \begin{eqnarray}
        &&\hat{\mathbb{G}}_{eff}^{(k)}\equiv \ep\: \mathcal{E}_g\lb \hat{\mathbb{H}}_g^2\rb=\sum_I
        \hat{\mathbb{X}}_I^2\nn\\
        &&\hat{\mathbb{X}}_I=\sqrt{\sigma}\sum_a \hat{V}_I^{(a)}=\mathbb{X}_{I,L}-\mathbb{X}_{I,R}^T\nn\\
        &&\mathbb{X}_{I}=\sqrt{\sigma}\sum_a V_I^{(a)}\ .
    \end{eqnarray}
    It follows that
    \begin{eqnarray}
        \mathcal{E}_g(dU^{\otimes k}\otimes (dU^*)^{\otimes k})= 1^{\otimes 2k}-\frac{\ep}{2}\hat{\mathbb{G}}_{eff}^{(k)}+O(\ep^2)\ .
    \end{eqnarray}
    Then, we find
    \begin{eqnarray}
        &&\mathcal{E}_g(dU^{\otimes k}\otimes (dU^*)^{\otimes k})(\mathbb{Y}\otimes
        1^{\otimes k})\ket{1}^{\otimes k}\nn\\
        &&=\lb 1^{\otimes 2k}-\frac{\ep}{2}\hat{\mathbb{G}}_{eff}^{(k)}\rb(\mathbb{Y}\otimes
        1^{\otimes k})\ket{1}^{\otimes k}\nn\\
        &&=\lb \lb 1-\ep \:\mathcal{G}_{eff}^{(k)}\rb(\mathbb{Y})\otimes 1^{\otimes 1}\rb\ket{1}^{\otimes k}
    \end{eqnarray}
    where $\mathbb{Y}$ is a general operator acting on $\mathcal{H}^{\otimes k}$,
    and the super-operator $\mathcal{G}$ generates a Lindblad evolution
    \begin{eqnarray}
        \mathcal{E}_g(\p_t \mathbb{Y})=-\mathcal{G}_{eff}^{(k)}(\mathbb{Y})&=&\sum_{I,a,a'}\lb
        V_I^{(a)} \mathbb{Y} V_I^{(a')}-\frac{1}{2}\{ V_I^{(a)}V_I^{(a')},\mathbb{Y}\}\rb\nn\\
        &=&\sum_I \lb \mathbb{X}_I Y \mathbb{X}_I-\frac{1}{2}\{\mathbb{X}_I^2,\mathbb{Y}\}\rb\nn\\
        &=&-\frac{1}{2}\sum_I [\mathbb{X}_I,[\mathbb{X}_I,\mathbb{Y}]]\ .
    \end{eqnarray}
    where we have used $[V_{I}^{(a)},V_{I}^{(a')}]=0$. Here, our Lindblad jump
    operators are the self-adjoint $\mathbb{X}_{I}$ Lindblad operators. Alternatively,
    we can write this Lindbladian in the Schrodinger picture for the density
    matrix by replacing $Y$ with $\rho^{\otimes k}$:
    \begin{eqnarray}
        \mathcal{E}_g(\p_t \rho^{\otimes k})=-\mathcal{G}_{eff}^{(k)}(\rho^{\otimes k})\ .
    \end{eqnarray}

    Integrating our Lindbladian evolution over a time-interval $(t,t+T)$ results
    in a Markov semi-group
    \begin{eqnarray}
        \mathcal{E}_{g;T}(\mathbb{Y})=e^{-T\mathcal{G}_{eff}^{(k)}}(\mathbb{Y})\ .
    \end{eqnarray}

    \section{Semigroup in Brownian SYK}

    The Brownian Sachdev-Ye-Kitaev (SYK) model is a variant of the standard SYK
    model where the static quenched disorder terms in the Hamiltonian are
    replaced by time-dependent stochastic variables \cite{saad2018semiclassical}.
    The Hamiltonian is given by:
    \begin{equation}
        \label{SYKHam}H(t) = i^{q/2}\sum_{i_1<\cdots <i_q}J_{i_1\cdots i_q}(t) (\chi
        _{i_1}\cdots \chi_{i_q})\ .
    \end{equation}

    The couplings $J_{I}(t)$ are independent Gaussian white noise processes defined
    by:
    \begin{align}
        \overline{J_I(t)}            & = 0, \nonumber                      \\
        \overline{J_I(t) J_{I'}(t')} & = \sigma\delta_{II'}\delta(t-t')\nn \\
        \sigma                       & =\frac{2^{q-1}q!}{q^{2}N^{q-1}}J\ .
    \end{align}

    Plugging (\ref{VSYK}) into the equation for $\hat{\mathbb{G}}_{eff}^{(k)}$,
    we obtain Lindblad operators
    \begin{eqnarray}
        &&\mathbb{X}_{i_1<\cdots <i_q}=\sum_a (\psi_{i_1}^{(a)}\cdots \psi_{i_q}^{(a)})=\sum_a
        (\chi_{i_1}^{(a)}\cdots \chi_{i_q}^{(a)})
    \end{eqnarray}
    and
    % The diagonal terms in the sum above, i.e. $a=b$, simplify to
    % \begin{align}
    % &\frac{2}{q!}\sum_{\text{distinct} \: i_1\cdots i_q}\sum_{a=1}^k\lb 1-\psi^{(aa)}_{i_1;LR}\cdots \psi^{(aa)}_{i_q;LR}\rb \nonumber\\
    % &=2\lb {N \choose q}-\sum_{i_1<\cdots<i_q}\sum_{a=1}^k\psi^{(aa)}_{i_1;LR}\cdots \psi^{(aa)}_{i_q;LR}\rb.\label{eq:diagonal_terms}
    % \end{align}
    % To compute the off-diagonal terms, i.e., $a\neq b$, we first note that

    % ----

    % Bilinear operators $\psi^{(ab)}_{ll';\alpha\beta}$ provide an orthonormal basis for the set of all operators in the full Hilbert space of $k N$ qubits:
    % \begin{eqnarray}
    % &&\psi^{(ab)}_{ll';\alpha\beta}\equiv \psi^{(a)}_{l \alpha}\psi^{(b)}_{l' \beta}\ .
    % \end{eqnarray}

    % %-----

    % Note that there is asymmetry between the $L$ and $R$ copies in the dictionary between the complex fermions and the qubits. Since we are ultimately interested in the Weyl fermions,
    % Note the asymmetry between $L$ and $R$.
    % As a result,
    % \begin{eqnarray}
    % &&\chi_{2k-1,L}=(-1)^{k-1}Q_{k-1}X_{kL},\qquad \chi_{2k,L}=(-1)^{k-1}Q_{k-1}Y_{kL}\nn\\
    % &&\chi_{2k-1,R}=(-1)^{k-1}Q_{k-1}(-Z_k)Y_k,\qquad \chi_{2k,R}=(-1)^{k-1}Q_{k-1}(-Z_k)X_k\nn\\
    % &&Z_{kL}=-i\chi_{2k-1,L}\chi_{2k,L},\qquad Z_{kR}=i\chi_{2k-1,R}\chi_{2k,R}\nn\\
    % &&Q_k=\prod_{j=1}^{2k}
    % \end{eqnarray}

    % Written in terms of the Weyl fermions, we find
    % \begin{eqnarray}
    % &&(\chi_{2k-1,L}+i\chi_{2k-1,R})\ket{e}^{\otimes N/2}=0=(\chi_{2k,L}+i\chi_{2k,R})\ket{e}^{\otimes N/2}\nn\\
    % && (\chi_{2k-1,L}-i\chi_{2k-1,R})\ket{f}^{\otimes N/2}=0=(\chi_{2k,L}+i\chi_{2k,R})\ket{f}^{\otimes N/2}\ .
    % \end{eqnarray}
    % Note that $(\chi_{2k,L}+i\chi_{2k-1,R})$ kills both vectors above. Notice that the vector $\ket{e}^{\otimes N/2}$ has the advantage that we do not need to distinguish between $\chi_{2k-1}$ and $\chi_{2k}$, we can simply write the mirror equation
    % \begin{eqnarray}
    % \chi_{j,L}\ket{e}^{\otimes N/2}=-i\chi_{j,R}\ket{e}^{\otimes N/2}\ .
    % \end{eqnarray}

    % In the GNS Hilbert space, the Hamiltonian of SYK for $q$ and $N$ both even can be written as
    % \begin{eqnarray}
    % \hat{H}=H_L-H_R=\sum_{i_1<\cdots<i_q}J_{i_1\cdots i_q}\lb \chi_{i_1L}\cdots \chi_{i_qL}-  \chi_{i_1R}\cdots \chi_{i_qR}\rb
    % \end{eqnarray}
    % Therefore, we can write $\hat{H}=H\otimes 1-1\otimes H$.

    % -----

    % We define the complex fermions
    % \begin{eqnarray}
    % \hat{\chi}_j^{(a)}=\chi^{(a),L}_j-i\chi^{(a),R}_j\ .
    % \end{eqnarray}
    % and we denote their corresponding vacuum with the infinite temperature TFD vector $\ket{1}$ is defined as
    % \begin{eqnarray}
    % \hat{\chi}^{(a)}_j\ket{1}^{(a)}=0\ .
    % \end{eqnarray}

    % Define the operators
    % \begin{eqnarray}
    % Q^{(a)}=\chi^{(a)}_1\cdots \chi^{(a)}_N
    % \end{eqnarray}
    % so that
    % \begin{eqnarray}
    % Q^{(a)}\chi_j^{(a)}=(-1)^{N-j}\chi^{(a)}_1\cdots \chi_{j-1}^{(a)}\chi^{(a)}_j\cdots \chi_N^{(a)}=(-1)^{N-1}\chi_j^{(a)}Q^{(a)}\ .
    % \end{eqnarray}
    % We assume that $N$ is even, so that
    % \begin{eqnarray}
    % &&Q^{(a)}\chi^{(j)}=-\chi_j^{(a)}Q^{(a)},\qquad Q^{(a)}Q^{(a)}=(-1)^{(N/2 \mod 2 )}\ .
    % \end{eqnarray}
    % Following \cite{}, we define
    % \begin{eqnarray}
    % \psi^{(a)}_j=i Q\chi^{(a)}_j
    % \end{eqnarray}
    % so that
    % \begin{eqnarray}
    % \chi^{(a)}_{i_1}\cdots \chi^{(a)}_{i_M}=(-1)^{N M}\psi^{(a)}_{i_1}\cdots \psi^{(a)}_{i_M}
    % \end{eqnarray}
    % where we are assuming both $N$ and $M$ are even.

    % For example, for the case $q=2$, the Lindblad operators are
    % \begin{eqnarray}
    % \mathbb{X}
    % \end{eqnarray}

    % We will see that thanks to the Brownian time-dependent noise, the averaged dynamical evolution of the system above is exactly solvable even at finite $N$. The time-evolution operator is
    % \begin{eqnarray}
    %     &&U(t)=\mathcal{T}\lb e^{-i\int_0^t dt' H(t')}\rb
    % \end{eqnarray}
    % and we are interested in computing
    % \begin{eqnarray}
    %     &&\overline{U(T)\otimes U(T)^\dagger}=\overline{\mathcal{T}\lb e^{-i\int_0^t dt\: \tilde{H}(t')}\rb}\nn\\
    %     &&\tilde{H}(t)=H_L(t)-H_R(t)=i^{q/2}\sum_I J_I(t)(\Psi_I^L-\Psi_I^R),
    % \end{eqnarray}
    % where $H_L$ and $H_R$ are two identical copies of $H$.

    % The Brownian averages over $J(t)$ can be done using a path-integral
    % \begin{eqnarray}
    %     &&\int \mathcal{D}J_I\: e^{-\frac{1}{2\sigma}\sum_I\int dt' J_I(t')^2}\ .
    % \end{eqnarray}
    % Therefore,
    % \begin{eqnarray}
    %     \overline{U(T)\otimes U(T)^\dagger}&&=\int\mathcal{D}J_I \exp\lb\frac{-1}{2\sigma}\sum_I \int_0^T dt'dt J_I(t')\delta(t'-t)\lb J_I(t)+2i\sigma \mathbb{X}_I\rb\rb\nn\\
    %     &&=e^{-\frac{\sigma T}{2}\sum_I \mathbb{X}_I^2}\int \mathcal{D}\tilde{J}_I\exp\lb \frac{-1}{2\sigma}\sum_I \tilde{J}_I^2\rb\nn\\
    %     &&=e^{-\frac{\sigma T}{2}\sum_I \mathbb{X}_I^2}\nn\\
    %     &&\tilde{J}=J+i\sigma \mathbb{X}_I\ .
    % \end{eqnarray}
    % In other words, we find
    % \begin{eqnarray}
    %    &&\overline{U(T)\otimes U(T)^\dagger}=e^{-T H_{eff}} \nn\\
    %    &&H_{eff}=\frac{\sigma}{2}\tilde{H}^2=\frac{\sigma(-1)^q}{2}\sum_I (\Psi^L_I-\Psi^R_I)^2\ .
    % \end{eqnarray}
    % This argument can be repeated for any $H=\sum_I J_I(t) H_I$. Since
    % \begin{eqnarray}
    %     \tilde{H} X\ket{1}=\sum_I [H_I,\mathcal{O}]\ket{1}\nn\ .
    % \end{eqnarray}
    % The super-operator that corresponds to $\tilde{H}_{eff}$ is
    % \begin{eqnarray}
    %     \sum_I [H_I,[H_I,\mO]]\ .
    % \end{eqnarray}

    \subsection{Time-dependent correlators}

    We can compute time-dependent correlators as well. The first observation is that
    we can always order times such that $t_{1}<t_{2}<\cdots <t_{k}$. Then, it follows
    from the calculation in the previous section that
    \begin{eqnarray}
        \mathcal{E}_g\lb\otimes_{i=1}^k \lb U(t_i)\otimes U(t_i)^*\rb\rb=e^{-(t_k-t_{k-1}) \mathbb{G}^{(1)}_{eff}}e^{-(t_2-t_1)\mathbb{G}^{(k-1)}_{eff}}e^{-t_1\mathbb{G}^{(k)}_{eff}}
    \end{eqnarray}
    where $G_{eff}^{(j)}$ is defined using (\ref{eq:Geff_sum_over_AB}) with $\alpha
    ,\beta=L,R$ but $a,b=1,\cdots j$.

    We would like to use this to compute correlation functions

    % \NL{What is the $k$-fold of the semigroup? expand in small $t$: $1-t\mathcal{L}$, and construct the probability measure over unitaries in this limit.}

    \section{Approaching \texorpdfstring{$k$}{k}-designs}

    We are interested in studying how quickly our Brownian dynamics approaches
    $k$-design. In other words, we would like to compare the $k$-th moment of
    our Brownian dynamics
    \begin{eqnarray}
        \mathcal{E}^{(k)}_{g;T}(\mathbb{Y})=e^{-T\mathcal{G}^{(k)}}(\mathbb{Y})
    \end{eqnarray}
    to that of a Haar random ensemble
    \begin{eqnarray}
        \mathcal{T}_H^{(k)}(\mathbb{Y})=\int d\mu_{Haar}(U)\:\lb U^{\otimes k}(\mathbb{Y})(U^\dagger)^{\otimes k}\rb
    \end{eqnarray}
    $d\mu(U)$ denotes the Haar measure, and we have added the superscript $(k)$ to
    keep track of the dimensionality of the space of operator $\mathbb{Y}\in \mathcal{L}
    ((\mathbb{C}^{d})^{\otimes k})$.

    \subsection{Haar Twirl}
    \label{subsec:haar_twirl}

    We start by reviewing the explicit form of the $k$-fold Haar twirl channel
    $\mathcal{T}^{(k)}$. By Schur-Weyl duality, the action of the unitary group
    on the tensor product space commutes with the action of the symmetric group
    $S_{k}$ permuting the tensor factors. As a result, the image of
    $\mathcal{T}^{(k)}$ is the algebra spanned by the permutation partial
    isometries $\{P_{\sigma}\}_{\sigma \in S_k}$:
    \begin{eqnarray}
        P_{\sigma}\ket{i_1\cdots i_k}=\ket{i_{\sigma(1)}\cdots i_{\sigma(k)}}\ .
    \end{eqnarray}
    Assuming $d \geq k$, these operators are linearly independent, allowing us to
    expand the twirled channel as
    \begin{eqnarray}
        \mathcal{T}^{(k)}(\mathbb{Y}) = \sum_{\sigma \in S_k} c_\sigma(\mathbb{Y})
        P_\sigma\ .
    \end{eqnarray}
    To determine the coefficients $c_{\sigma}$, we exploit the fact that
    $\mathcal{T}^{(k)}$ is an orthogonal projection under the Hilbert-Schmidt
    inner product, and
    \begin{eqnarray}
        \Tr(P_\tau^\dagger \mathcal{T}^{(k)}(\mathbb{Y}))&&=\int d\mu_{Haar}(U)\Tr\lb
        (U^\dagger)^{\otimes k} P_\tau^\dagger U^{\otimes k}\: \mathbb{Y}\rb\nn\\
        &&=\int d\mu_{Haar}(U)\Tr\lb P_\tau^\dagger \: \mathbb{Y}\rb=\Tr\lb P_\tau^\dagger
        \mathbb{Y}\rb
    \end{eqnarray}
    where we have used $P_{\tau}^{\dagger}U^{\otimes k}=U^{\otimes k}P_{\tau}^{\dagger}$.
    Substituting the expansion for $\mathcal{T}^{(k)}(X)$ yields a linear system
    involving the Gram matrix of the permutation basis, $G_{\tau, \sigma}= \Tr(P_{\tau}
    ^{\dagger}P_{\sigma}) = d^{\#(\tau^{-1}\sigma)}$, where $\#(\pi)$ counts the
    disjoint cycles in $\pi$. Inverting this relation requires the unitary Weingarten
    function, defined as $Wg(\pi, d) = [G^{-1}]_{\pi}$. This yields the explicit
    closed-form expression for the $k$-fold Haar twirl channel:
    \begin{equation}
        \mathcal{T}^{(k)}(\mathbb{Y}) = \sum_{\sigma, \tau \in S_k}Wg(\sigma\tau^{-1}
        , d) \Tr(P_{\tau}^{\dagger}\mathbb{Y}) P_{\sigma}. \label{eq:general_twirl}
    \end{equation}
    The Weingarten functions are discussed extensively in \cite{collins2006}. As
    a concrete example, for the $k=2$ case relevant to standard randomized
    benchmarking, the symmetric group is $S_{2}= \{e, (12)\}$. The Weingarten
    matrix is the inverse of the Gram matrix $\begin{pmatrix}
        d^{2} & d     \\
        d     & d^{2}
    \end{pmatrix}$, given by
    \begin{equation}
        Wg = \frac{1}{d^{2}- 1}
        \begin{pmatrix}
            1       & -d^{-1} \\
            -d^{-1} & 1
        \end{pmatrix}.
    \end{equation}
    Substituting these values into Eq.~\eqref{eq:general_twirl}, we recover the
    standard projection onto the isotropic (Werner) states:
    \begin{equation}
        \mathcal{T}^{(2)}(\mathbb{Y}) = \frac{\Tr(\mathbb{Y}) - d^{-1}\Tr(\mathbb{Y}S)}{d^{2}-1}
        \mathbb{I}+ \frac{\Tr(\mathbb{Y}S) - d^{-1}\Tr(\mathbb{Y})}{d^{2}-1}S.
    \end{equation}
    where $S$ is the swap operator.

    \subsection{Distance measures}

    To compare our Brownian dynamics to the Haar averages, we use three different
    distance measures:
    \begin{enumerate}
        \item {\bf Hilbert-Schmidt (Frobenius) distance:} This measure is the $(2
            \to 2)$-norm defined of the difference of our superoperators as
            \begin{eqnarray}
                \delta^{HS}_k(T)\equiv \|\mathcal{E}_{g;T}^{(k)}-\mT^{(\otimes)}_H\|_{2\to 2}=\sup_{\mathbb{Y}\neq 0}
                \frac{\|\mathcal{E}_{g;T}^{(k)}(\mathbb{Y})-\mT^{(\otimes)}_{H}(\mathbb{Y})\|_{2}}{\|\mathbb{Y}\|_{2}}
            \end{eqnarray}
            where $\|\mathbb{Y}\|_{2}=\sqrt{\Tr(\mathbb{Y}^{\dagger}\mathbb{Y})}$.
            We will see that for our semigroup, the existence of a gap $\gamma>0$
            in the spectrum of the generator implies that the Frobenius norm decays
            exponentially in time $\gamma T$.

        \item {\bf Diamond distance:} This is the most standard distance measure
            in comparing finite dimensional quantum channels:
            \begin{eqnarray}
                \delta_k^{\diamond}(T)\equiv\sup_m\|(\mathcal{E}_{g;T}^{(k)}-\mT^{(\otimes)}_H)\otimes
                1_m\|_{1\to 1}=\sup_{\mathbb{Y}>0,m}
                \frac{\|((\mathcal{E}_{g;T}^{(k)}-\mT^{(\otimes)}_{H})\otimes 1_{m})(\mathbb{Y})\|_{1}}{\|\mathbb{Y}\|_{1}}
            \end{eqnarray}

        \item {\bf Relative entropy distance:} The relative entropy distance is defined
            as
            \begin{eqnarray}
                \delta^{RE}_k(T)\equiv\sup_m\sup_\rho D((\mathcal{E}^{(k)}_{g;T}\otimes
                1_m)(\rho)\|(\mathcal{T}^{(k)}_H\otimes 1_m)(\rho))
            \end{eqnarray}
            This type of distance measure is often studied using logarithmic Sobolev
            inequalities.

        \item {\bf $\ep$-Relative errors:} The strongest distance measure we consider
            is
            \begin{eqnarray}
                \ep_k(T)\equiv \inf_\ep\{\ep>0:(1-\ep)(\mathcal{E}^{(k)}_{g;T}\circ
                \mathcal{T}^{(k)}_H)\leq \mathcal{E}^{(k)}_H\leq (1+\ep)(\mathcal{E}^{(k)}_{g;T}\circ
                \mathcal{T}^{(k)}_H)\:\}
            \end{eqnarray}
            If all permutations remain invariant under $\mathcal{E}^{(k)}_{g;T}$,
            then our distance measure simplifies to
            \begin{eqnarray}
                \ep_k(T)\equiv \inf_\ep\{\ep>0:(1-\ep)\mathcal{T}^{(k)}_H\leq \mathcal{E}^{(k)}_H\leq
                (1+\ep)\mathcal{T}^{(k)}_H\:\}
            \end{eqnarray}
    \end{enumerate}
    Note that a priori, it is not clear that any of the distance measures above has
    a large $T$ limit in our dynamical system of interest.

    \subsection{Steady states of Brownian dynamics}

    The steady states of our Brownian dynamics are those that are killed by the
    Lindbladian generator $\mathcal{G}_{eff}^{(k)}$.
    \begin{equation}
        \mathcal{G}_{eff}^{(k)}(\mathbb{Y}) = -\sum_{I}\left( \mathbb{X}_{I}\mathbb{Y}
        \mathbb{X}_{I}- \frac{1}{2}\{ \mathbb{X}_{I}^{2}, \mathbb{Y}\} \right) =
        \frac{1}{2}\sum_{I}[\mathbb{X}_{I}, [\mathbb{X}_{I}, \mathbb{Y}]]. \label{eq:generator_form}
    \end{equation}
    For simplicity, we have absorbed the variance $\sigma$ (cf.\ its definition
    in the Brownian ensemble) into the definition of $\mathbb{X}_{I}$. We define
    the \textit{spectral gap} $\Delta_{k}$ as the smallest real part of the non-zero
    eigenvalues of $\mathcal{G}^{(k)}$. The following theorem governs the
    convergence of the dynamics to the Haar average.

    \begin{theorem}[Convergence to steady states]
        Let $\mathfrak{L}\subseteq \mathfrak{su}(d)$ be the Lie algebra generated
        by the set of operators $\{ V_{I}\}_{I}$. The following three statements
        are equivalent:
        \begin{enumerate}
            \item \textbf{Algebraic universality:} The algebra $\mathfrak{L}$
                acts irreducibly on $\mathcal{H}$; i.e.
                $\mathfrak{L}= \mathfrak{u}(d)$.

            \item \textbf{Steady state space:} The kernel of the generator
                coincides exactly with the commutant of the unitary group, i.e.,
                $\ker(\mathcal{G}^{(k)}) = \mathrm{span}\{P_{\sigma}\mid \sigma \in
                S_{k}\}$.

            \item \textbf{Strictly Positive Gap:} The spectral gap satisfies
                $\Delta_{k}> 0$, implying exponential convergence to the Haar twirl
                in any norm:
                \begin{equation}
                    \| e^{-T\mathcal{G}^{(k)}}- \mathcal{T}^{(k)}\|_{2\to2}\leq e
                    ^{-\Delta^{(k)} T}.
                \end{equation}
                Since we are working with finite-dimensional matrix algebras, this
                implies convergence in the diamond distance and relative entropy
                distance, as well:
                \begin{eqnarray}
                    \lim_{T\to \infty}\delta^{\diamond}_k(T)=0=\lim_{T\to \infty}\delta^{RE}_k(T)
                \end{eqnarray}
        \end{enumerate}
    \end{theorem}

    \begin{proof}
        We proceed by establishing the chain of implications $(1) \implies (2) \implies
        (3) \implies (1)$.

        \paragraph{$(1) \implies (2)$:}
        Viewing $\mathcal{L}(\mathcal{H}^{\otimes k})$ as a Hilbert space with
        the Hilbert-Schmidt inner product
        $\langle A, B \rangle = \Tr(A^{\dagger}B)$, we have
        \begin{equation}
            \langle \mathbb{Y}, \mathcal{G}^{(k)}(\mathbb{Y}) \rangle = -\frac{1}{2}
            \sum_{I}\| [\mathbb{X}_{I}, \mathbb{Y}] \|_{2}^{2}.
        \end{equation}
        Since norms are non-negative, $\mathcal{G}^{(k)}(\mathbb{Y}) = 0$ if and
        only if $[\mathbb{X}_{I}, \mathbb{Y}] = 0$ for all $I$. This implies that
        $\mathbb{Y}$ must commute with the entire Lie algebra generated by the global
        operators $\{\mathbb{X}_{I}\}$. Since
        $\mathbb{X}_{I}=\sqrt{\sigma}\sum_{a=1}^{k}V_{I}^{(a)}$, this condition is
        equivalent to commuting with the group $G^{\otimes k}$ generated by exponentiating
        the algebra $\mathfrak{L}$. If $\mathfrak{L}= \mathfrak{su}(d)$, then $G
        = SU(d)$. By Schur-Weyl duality, the operators acting on $(\mathbb{C}^{d}
        )^{\otimes k}$ that commute with $U^{\otimes k}$ for all $U \in SU(d)$ are
        precisely the linear combinations of permutation operators $P_{\sigma}$.
        Thus, $\ker(\mathcal{G}^{(k)}) = \text{span}(S_{k})$.

        \paragraph{$(2) \implies (3)$:}
        The superoperator $\mathcal{G}^{(k)}$ viewed as an operator on the GNS
        Hilbert space is $\hat{\mathbb{G}}^{(k)}$
        \begin{eqnarray}
            &&\mathcal{G}^{(k)}(\mathbb{Y})\ket{1}^{\otimes k}=\hat{\mathbb{G}}^{(k)}(\mathbb{Y}\otimes
            1^{\otimes k})\ket{1}^{\otimes k}\nn\\
            &&\hat{\mathbb{G}}^{(k)}=\sum_I \hat{\mathbb{X}}_I^2
        \end{eqnarray}
        which is manifestly positive semi-definite. Therefore, the spectrum of
        the superoperator $\mathcal{G}^{(k)}$ is comprised of non-negative
        eigenvalues $\lambda$'s. Statement (2) implies that the zero eigenspace ($\lambda
        =0$) corresponds exactly to the image of the Haar twirl $\mathcal{T}^{(k)}$.
        Since the dimension is finite, the spectrum is discrete. Therefore, the smallest
        non-zero eigenvalue, defined as the gap
        $\Delta^{(k)}= \min \{ \lambda_{j}\mid \lambda_{j}> 0 \}$, is strictly positive.
        Standard spectral theory implies that the evolution on the orthogonal
        complement of the kernel decays as
        $\| e^{-T\mathcal{G}^{(k)}}|_{\ker^\perp}\| \le e^{-\Delta^{(k)} T}$.

        \paragraph{$(3) \implies (1)$:}
        We prove by contradiction. Suppose (3) statement holds but (1) does not
        hold; i.e., the generated Lie algebra $\mathfrak{L}$ is a proper subalgebra
        of $\mathfrak{su}(d)$. The dynamical group $G = e^{\mathfrak{L}}$ is
        then a proper subgroup of $SU(d)$. The commutant of a subgroup is strictly
        larger than the commutant of the full group (except in trivial cases, which
        are excluded by the assumption of proper subalgebra). That is,
        $\dim(\ker(\mathcal{G}^{(k)})) = \dim(\text{Comm}(G)) > \dim(\text{span}(
        S_{k}))$. This implies there exists at least one operator $O \in \ker(\mathcal{G}
        ^{(k)})$ such that $O \notin \text{span}(S_{k})$. Consequently, the steady
        state of the evolution includes $O$. The distance to the Haar twirl
        $\mathcal{T}^{(k)}$ (which projects out $O$) will never decay to zero, implying
        the effective gap for convergence to $\mathcal{T}^{(k)}$ is zero.
    \end{proof}
    If the Lie algebra generated by $\mathbb{X}_{I}$ is a proper subalgebra of $\mathfrak{su}
    (d)$, the missing generators generate constants of motion, and should be physically
    interpreted as conserved charges.

    \section{Conserved charges in Brownian SYK}

    We would like compute the commutant of the Brownian $k$-twirl generators. Let
    \[
        \mathrm{Comm}\{\mathbb{X}_{I}\}\;:=\;\Big\{Y\in\mathrm{End}(\mathcal{H}^{\otimes
        k})\;:\;[\mathbb{X}_{I},Y]=0\ \ \forall I\Big\}, \qquad \mathbb{X}_{I}=\sqrt{\sigma}
        \sum_{a=1}^{k}V_{I}^{(a)}.
    \]
    Since different replicas commute, the Lie algebra generated by
    $\{\mathbb{X}_{I}\}$ is the diagonal action of the single-copy dynamical
    group $G:=\langle e^{itV_I}\rangle$, i.e.
    \[
        \mathrm{Comm}\{\mathbb{X}_{I}\}=\{Y:\ [Y,U^{\otimes k}]=0\ \forall U\in G
        \},
    \]
    the centralizer algebra of the diagonal representation.

    \paragraph{Explicit form for Brownian SYK (even $q$).}
    For SYK with even $q$, each
    $V_{I}^{(a)}=i^{q/2}\chi^{(a)}_{i_1}\cdots\chi^{(a)}_{i_q}$ (equivalently with
    $\psi$ in place of $\chi$) is an even fermion operator and commutes with the
    parity on that replica,
    \[
        Q^{(a)}:=(-i)^{N/2}\prod_{j=1}^{N}\chi^{(a)}_{j}\quad (\text{or }(-i)^{N/2}
        \prod_{j}\psi^{(a)}_{j}),\qquad [Q^{(a)},\mathbb{X}_{I}]=0\ \ \forall a,
        I.
    \]
    Moreover replica permutations $P_{\sigma}$ commute with $\mathbb{X}_{I}$ because
    they relabel the replica sum:
    $P_{\sigma}\mathbb{X}_{I}P_{\sigma}^{-1}=\mathbb{X}_{I}$. Hence always
    \[
        \mathrm{Alg}\langle\,P_{\sigma},\ Q^{(1)},\dots,Q^{(k)}\,\rangle\ \subseteq
        \ \mathrm{Comm}\{\mathbb{X}_{I}\}.
    \]
    If the SYK generators $\{V_{I}\}$ act irreducibly on each fixed-parity sector
    (equivalently, generate the full even operator algebra on one copy), then this
    inclusion is an equality and a convenient linear basis is
    \[
        \boxed{\quad \mathrm{Comm}\{\mathbb{X_I}\} =\mathrm{span}\Big\{\,P_\sigma\,\prod_{a=1}^k (Q^{(a)})^{\epsilon_a}\ :\ \sigma\in S_k,\ \epsilon_a\in\{0,1\}\Big\}. \quad}
    \]
    Equivalently, using parity projectors
    $\Pi_{\vec s}:=\prod_{a=1}^{k}\frac{1+s_{a}Q^{(a)}}{2}$ with $s_{a}=\pm1$, one
    may take the basis $\{\Pi_{\vec s}P_{\sigma}\}$. In the common setting where
    each copy is restricted to a fixed parity sector, $Q^{(a)}=\pm1$ is scalar and
    the commutant reduces to $\mathrm{span}\{P_{\sigma}\}$.

    \subsection{Old text}

    \NL{Use path-integrals to write down the time-evolved operator above as a Euclidean path-integral}
    To start, consider the quantum mechanics of a single particle, where the position
    of a particle $q_{I}(t)$ in $\mathbb{R}^{n}$ with $I=1,\cdots, n$ is an operator\footnote{In
    QFT language, $q_{I}(t)$ is $n$ copies of $0+1$-dimensional quantum fields.}.
    For a particular history of couplings $g(t)$, we can represent the
    correlator above using a path-integral

    \NL{Nick: we would like to do the whole analysis directly in the continuum, and not bring up brickworks.}

    \NL{Is there an ergodicity statement in the $k$-fold establishing that the $k$-fold semigroup generates the whole unitary ensemble. This is equivalent to the absence of other charges}

    \NL{For actual SYK, we expect a semi-group to emerge at very low frequency $(\beta\omega)\ll (\beta g)$. is there something related to this that we can say? for example, can we replace the Brownian noise with time-independent noise, and use a low-pass filter to project to low frequencies...}

    \section{To do}

    The form of the generators of the Lindbladian is a sum over a hypergraph with hyperedges indexed by $\{i_1,\cdots ,i_q\}$ distinct. go to eq 9 of Nick's notes, and generalize to hypergraphs, with a fermion on each site. 

    Work to bring eq 9 closer to the SYK case by putting fermions on each vertex. 

    Does the gluing lemma hold when you put fermionic degrees of freedom on vertices?

    
    
    
    Then, write the term inside in terms of $k$ qubits (by Jordan Wigner) and compare to eq 9 of Nick.

    \bibliographystyle{unsrt}
    \bibliography{ref}
\end{document}